\documentclass{zhvt-classic}

% 使用本地源流明体（GenRyuMin2 TW）按文件路径加载，不依赖系统字体名
% 版框（内框／界栏／书口侧栏／鱼尾）改用与版心印章同一印泥棕色：
%   印章图片 yshs.png 实色 = RGB(135,68,52) = HTML 874434
\usepackage{xcolor}
\usepackage{tikz}
% xcolor 与 l3color 是两套色彩注册表，需各注册同名 sealbrown：
%   xcolor 版 → 供 TikZ/pgf 的 fill= 读取；
%   L3 版     → 供 cls 版框 \color_stroke:n 读取。
\definecolor{sealbrown}{HTML}{874434}
\ExplSyntaxOn
\color_set:nnn { sealbrown } { HTML } { 874434 }
\ExplSyntaxOff

\zhvtset{
  grid color       = sealbrown,
  main font        = GenRyuMin2TW-SB,
  main font path   = C:/Users/Administrator/texmf/fonts/opentype/public/GenRyuMin2TW/,
jiazhu font      = GenRyuMin2TW-M,
  jiazhu font path = C:/Users/Administrator/texmf/fonts/opentype/public/GenRyuMin2TW/,
  % 页码框容纳 4 字：正常页「九十七」+前置空 1 格=4 行恰好；
  % 100+ 页「一百零一」共 5 行超出，末字「一/二/三」被框裁掉（2026-10-04 实测），故放宽到 5
  shukou page number chars = 6,
  % 书口书名框高从默认 6 字减到 4 字，把省出的 2 字高拨给 chapter 框，
  % 使 6 字书名+卷名放进 9 字高的 chapter 框后不再压线（2026-10-05）
  shukou title chars = 4,
}

\title{史记}[八书]
\author{汉·司马迁}
\maker{闲叟漫记·重製}

% ★ 修正书口卷名位置：chapter 框起始位置改回 cls 默认值，避免全局影响。
%   只在每个卷首页通过 \juan 命令里的空命令微调该页位置（安全，不会跨页污染）。
\ExplSyntaxOn
\AtBeginDocument{
  \dim_set:Nn \l__zhvt_shukou_chapter_title_start_dim
    { \l__zhvt_shukou_title_height_dim + \l__zhvt_font_size_dim * 2 }
  \dim_set:Nn \l__zhvt_shukou_page_number_start_dim
    { \l__zhvt_shukou_chapter_title_start_dim + \l__zhvt_shukou_chapter_title_height_dim }
}
\ExplSyntaxOff

% ★ 述赞样式：题头“索隐述赞”四字用黑体（simhei，与正文宋体/仿宋区分）+ 圆角棕框反白；
%   正文韵文按最初要求用仿宋（simfang），整体缩小字号（21pt）以便每四句单独成列、居列中。
%   ⚠ 竖排关键：自定义字体家族必须带 RawFeature={vertical:+vert}（仅 +vert）。
%      simhei 若加 +vhal 会整体加载失败回退成仿宋；simfang 同理只用 +vert。
\newCJKfontfamily{\shuzanfont}[Path=C:/Windows/Fonts/,RawFeature={vertical:+vert}]{simhei.ttf}
\newCJKfontfamily{\shuzanbody}[Path=C:/Windows/Fonts/,RawFeature={vertical:+vert}]{simfang.ttf}
% 圆角方框题头（2026-11-08 三次修正 · 定稿：字不得越框、不得有重影、墨迹真居中）：
%   ★ 弃用「隐形棕字节点 + 偏移 fill + 白字」三步法：text=sealbrown 的棕字层并非真隐形，
%     露出框外的部分会在白底上显出淡影，视觉上就是「反白字超出底纹」。
%   ★ 改用「定尺寸框 + 单层白字」：框 132pt × 38pt（绝对坐标），圆角 4pt。
%   ★★ 最大陷阱：pymupdf 的 span bbox 是 **advance 字身盒**，而 simhei 竖排替换字形
%      (vert) 的**实际墨迹相对字身盒整体左偏约 18pt**。只看 bbox 会误判「字在框内」，
%      实际笔画已被框左边切断。验收必须做像素级墨迹测量（渲染后统计框内白像素）。
%   ★ 校准结论（像素实测）：白字节点置于 TikZ (0pt, 8.2pt) —— x 方向**不需补偿**
%      （墨迹水平本就居中，早先的 -19.6/-9.6 是拿字身盒当墨迹造成的过补偿）；
%      y 方向 +8.2pt 抵消竖排墨迹垂直偏移。实测墨迹心与框心 dx=0.0 dy=0.0，
%      四边留白 L/R≈21pt、T/B≈8.5pt，且 21pt > 圆角 4pt，绝不切字越框。
%   页面 /Rotate 90 后，横长框自然成为竖条、四字竖排一列。
\newcommand{\shuzanhead}[1]{%
  \begin{tikzpicture}[baseline={(0,2)}]
    \fill[sealbrown, rounded corners=6pt]
      (-76pt,-15pt) rectangle (76pt,15pt);
    \node[text=white, inner sep=0pt, anchor=center] at (0pt,8.2pt)
      {\shuzanfont\fontsize{23}{23}\selectfont #1};
  \end{tikzpicture}%
}


\ExplSyntaxOn
\tl_set:Nn \l__zhvt_chapter_title_tl { 八书 }
\ExplSyntaxOff
\RenewDocumentCommand \tableofcontents { }
  {
    \clearpage
    史记八书目录
\par\hspace{3.5em}八书　凡八篇
\par\hspace{3.5em}礼书　　　　　\hspace{1.5em}卷二十三
\par\hspace{3.5em}乐书　　　　　\hspace{1.5em}卷二十四
\par\hspace{3.5em}律书　　　　　\hspace{1.5em}卷二十五
\par\hspace{3.5em}历书　　　　　\hspace{1.5em}卷二十六
\par\hspace{3.5em}天官书　　　　\hspace{1.5em}卷二十七
\par\hspace{3.5em}封禅书　　　　\hspace{1.5em}卷二十八
\par\hspace{3.5em}河渠书　　　　\hspace{1.5em}卷二十九
\par\hspace{3.5em}平准书　　　　\hspace{1.5em}卷三十
\par\hspace{3.5em}凡八篇　卷二十三至卷三十
    \clearpage
  }

\begin{document}

% ★★ 版心（书口）排布覆写：把书名放到鱼尾之下（2026-10-04）
%
%   cls 默认版心自上而下：书名（title 静态框，在鱼尾之上）→ 鱼尾 → 篇名（chapter 动态框）→ 页码
%   本模板改为：          鱼尾（位置不动）→ 书名 → 篇名 → 页码分割线 → 页码
%
%   手法：把 cls 放在 title 静态框里的书名，改放进 chapter 动态框（鱼尾之下）；
%         title 静态框留空。书名取 \l__zhvt_title_str，即 \title{...} 的第一参数，
%         所以换书时只改 \title，不必在此重复写书名。
%
%   ⚠ 三条必须遵守的写法（2026-10-04 实测）：
%     ① 汉字必须用 \hbox{...} 包住。flowfram 会把框内容装进一个宽 14.4mm（= 一个 line_sep
%        = 书口栏宽）的纵向 minipage 并按该宽度折行，而整页靠 PDF /Rotate 90 显示、
%        框内文字本身不变向 —— 不包 \hbox 时「围炉夜话」4 字会被逐字折成 4 行，
%        在版面上就成了横排 4 字并横向溢出版框。
%     ② 竖排下**一个 \hbox 就是一整列**。书名与篇名要同处一列，必须放进同一个 \hbox；
%        分成两个 \hbox 时，后一个会被当成新的一列、甩到隔壁栏去
%        （实测「序」跑到书口栏外侧 14.4mm 的正文列上）。中间的 \quad 在框内是横向空白，
%        显示出来正好是书名与篇名之间的纵向间隔。
%     ③ ★ 内容盒要写成 **与框等宽的定宽盒** `\hbox to \l__zhvt_shukou_chapter_title_height_dim{…\hfil}`。
%        实测：chapter 框是把内容盒沿纵向**居中**摆放的，内容宽度一变，书名就跟着漂：
%        自然宽度的 \hbox（书名 38.4mm）落在 x=157.2mm，加篇名后（57.6mm）落在 x=152.4mm，
%        目录页（无篇名）与正文页的书名就差半格。做成等宽盒后，无论有没有篇名，
%        内容固定从框顶起排，实测两页均落在 x=142.8mm，不再漂。
\ExplSyntaxOn
% 把书口 chapter 框（书名+篇名）起始位置下移 1 字高，避免 4 字卷名顶到鱼尾
% ⚠ cls 的 \__zhvt_calculate_dimensions: 在 \AtBeginDocument 里执行，会重算该 dim；
%   必须把覆写挂到 begin document 之后，否则会被二次覆盖（2026-10-05 实测）。
% (旧方案已移到导言区，直接修改 chapter 框起始位置)
% 卷标题命令 #1=书口篇名(鱼尾) #2=正文卷名（含卷序，卷名内空格用 ~）
%
% ★★ 目录页两项修订（2026-10-05，愚大人指出）★★
%   ① 页码起点统一对齐：
%      原状 —— 目录项写成 \par\hspace{4em}卷名\hspace{4em}页码，页码起点 = 卷名实宽 + 4em。
%      卷名实宽随字数变：「卷一~梁惠王上」6 字 +~、「卷七~离娄上」5 字 +~，
%      差 1 个字位 → 卷七~卷十 4 条的页码整体比其余 10 条高 27.21pt（一个字），看着上下不齐。
%      改法 —— 卷名补全角空格（宽 = 1 字 = 27.215pt）补足到 **7 字符位**再写进目录，
%      14 条页码起点遂全部落在同一条竖线上（起点对齐；末字随页码字数自然参差）。
%      ⚠ 只对目录里的副本 \l_tmpc_tl 补，正文卷名仍用 \l_tmpa_tl 原样 —— 正文像素不受影响。
%   ② 卷名整体上移半字位：前导 \hspace 由 4em 改 3.5em。
%      实测原状目录块在版框内偏下（天头留白 139.6pt / 地脚留白 110.7pt），
%      上移 0.5em = 13.6pt 后上下基本均分，视觉居中。
\NewDocumentCommand { \juan } { mm } {
	\clearpage
	\tl_set:Nn \l__zhvt_chapter_title_tl { #1 }
	\tl_set:Nn \l_tmpa_tl { #2 }
	\tl_set_eq:NN \l_tmpc_tl \l_tmpa_tl
	\int_step_inline:nn { 7 - \str_count:n { #2 } }
	  { \tl_put_right:Nn \l_tmpc_tl { 　 } }
	\tl_set:Nx \l_tmpb_tl { \zhnum { page } }
	\tl_put_left:Nn \l_tmpb_tl { \hspace { 4em } }
	\tl_put_left:Nn \l_tmpb_tl { \l_tmpc_tl }
	\tl_put_left:Nn \l_tmpb_tl { \par \hspace { 3.5em } }
	\addtocontents { toc } { \l_tmpb_tl }
	\skip_horizontal:n { \l__zhvt_font_size_dim * 4 }
	\l_tmpa_tl
	\par
	\thispagestyle { empty }
}
\cs_set_protected:Npn \__zhvt_set_shukou_content: {
	% ★ 鱼尾上方留白（title 静态框）嵌入闲叟钤印：
	%   整页 /Rotate 90，故此处把图用 \rotatebox{90} 转正为竖章；
	%   尺寸卡在版心条幅内（书内约 7mm 宽 × 33.8mm 高，留余量不压鱼尾）。
	\setstaticcontents * { title } {
		\vbox_to_ht:nn { \l__zhvt_baseline_skip_dim } {
			\vskip 14.2pt
			\hbox_to_wd:nn { \l__zhvt_shukou_title_height_dim } {
				\hfil
				\rotatebox [ origin = c ] { 90 } { \includegraphics [ width = 7mm , height = 33.8mm ] { yshs.png } }
				\hfil
			}
		}
	}
% ★ 版心（书口）统一改法（2026-11-16 愚大人定）：18pt＋框内居中＋书口文字同版框色
%   ① 原 27.2pt 时 7 字长篇名（如「廉颇蔺相如列传」）末字溢出 chapter 框、压住页码；
%      改 \fontsize{18pt}{18pt}，前后 \hfil 令字串在框内居中（不随篇名字数漂移）。
%   ② \vbox{\vskip 8.2pt} 定量下移，抵消缩字号后栏宽方向 8.2pt 上偏（像素实测 y 中心 543.3 = 原版）。
%      ⚠ \raisebox / \vspace* 在 flowfram 框内无效；换字号须同步换补偿：16pt→10.1、19pt→7.4。
%   ③ chapter 与 page~number 均置 \color{sealbrown}（HTML 874434），与版框／界行同色；正文仍黑。
%   ⚠ 页码字位、钤印、鱼尾、界行、正文一律未动；五部合订本同此口径。
	\setdynamiccontents * { chapter }
	{ \color{sealbrown} \vbox { \vskip 8.2pt \hbox to \l__zhvt_shukou_chapter_title_height_dim
		{ \hfil \fontsize{18pt}{18pt}\selectfont \l__zhvt_title_str \quad \l__zhvt_chapter_title_tl \hfil } } }
	\setdynamiccontents * { page~number } { \color{sealbrown} \quad \zhnum { page } }
}
\__zhvt_set_shukou_content:
\ExplSyntaxOff


\gridall
\tableofcontents

% ===== 卷二十三　礼书　书口/扉页篇名=礼书 =====
\clearpage
\ExplSyntaxOn
\tl_set:Nn \l__zhvt_chapter_title_tl { 礼书 }
\tl_set:Nn \l__zhvt_title_display_tl { 礼书 }
\ExplSyntaxOff
\maketitle

% 十表八书补排：单篇成册，不渲染目录页




% 方案①：仅改 A 类错讹，头部与版式不变

% ===== 正文（八书续排：vRain 白文底本 023，W 盘中华本核对，造字依定稿口径回改）=====
\juan{礼书}{卷二十三~礼书第一}
\hspace*{2em}太史公曰。洋洋美德乎。宰制万物。役使群众。岂人力也哉。余至大行礼官。观三代损益。乃知缘人情而制礼。依人性而作仪。其所由来尚矣。\\
\hspace*{2em}人道经纬万端。规矩无所不贯。诱进以仁义。束缚以刑罚。故德厚者位尊。禄重者宠荣。所以总一海内而整齐万民也。人体安驾乘。为之金舆错衡以繁其饰。目好五色。为之黼黻文章以表其能。耳乐钟磬。为之调谐八音以荡其心。口甘五味。为之庶羞酸咸以致其美。情好珍善。为之琢磨圭璧以通其意。故大路越席。皮弁布裳。朱弦洞越。大羹玄酒。所以防其淫侈。救其雕敝。是以君臣朝廷尊卑贵贱之序。下及黎庶车舆衣服宫室饮食嫁娶丧祭之分。事有宜适。物有节文。仲尼曰。﹁禘自既灌而往者。吾不欲观之矣。﹂。\\
\hspace*{2em}周衰。礼废乐坏。大小相逾。管仲之家。兼备三归。循法守正者见侮于世。奢溢僭差者谓之显荣。自子夏。门人之高弟也。犹云﹁出见纷华盛丽而说。入闻夫子之道而乐。二者心战。未能自决﹂。而况中庸以下。渐渍于失教。被服于成俗乎。孔子曰﹁必也正名﹂。于卫所居不合。仲尼没后。受业之徒沈湮而不举。或适齐。楚。或入河海。岂不痛哉。\\
\hspace*{2em}至秦有天下。悉内六国礼仪。采择其善。虽不合圣制。其尊君抑臣。朝廷济济。依古以来。至于高祖。光有四海。叔孙通颇有所增益减损。大抵皆袭秦故。自天子称号下至佐僚及宫室官名。少所变改。孝文即位。有司议欲定仪礼。孝文好道家之学。以为繁礼饰貌。无益于治。躬化谓何耳。故罢去之。孝景时。御史大夫晁错明于世务刑名。数干谏孝景曰。﹁诸侯藩辅。臣子一例。古今之制也。今大国专治异政。不禀京师。恐不可传后。﹂孝景用其计。而六国畔逆。以错首名。天子诛错以解难。事在袁盎语中。是后官者养交安禄而已。莫敢复议。\\
\hspace*{2em}今上即位。招致儒术之士。令共定仪。十余年不就。或言古者太平。万民和喜。瑞应辨至。乃采风俗。定制作。上闻之。制诏御史曰。﹁盖受命而王。各有所由兴。殊路而同归。谓因民而作。追俗为制也。议者咸称太古。百姓何望。汉亦一家之事。典法不传。谓子孙何。化隆者闳博。治浅者褊狭。可不勉与。﹂乃以太初之元改正朔。易服色。封太山。定宗庙百官之仪。以为典常。垂之于后云。\\
\hspace*{2em}礼由人起。人生有欲。欲而不得则不能无忿。忿而无度量则争。争则乱。先王恶其乱。故制礼义以养人之欲。给人之求。使欲不穷于物。物不屈于欲。二者相待而长。是礼之所起也。故礼者养也。稻粱五味。所以养口也。椒兰芬茝。所以养鼻也。钟鼓管弦。所以养耳也。刻镂文章。所以养目也。疏房床笫几席。所以养体也。故礼者养也。\\
\hspace*{2em}君子既得其养。又好其辨也。所谓辨者。贵贱有等。长少有差。贫富轻重皆有称也。故天子大路越席。所以养体也。侧载臭茝。所以养鼻也。前有错衡。所以养目也。和鸾之声。步中武象。骤中韶濩。所以养耳也。龙旗九斿。所以养信也。寝兕持虎。鲛韅弥龙。所以养威也。故大路之马。必信至教顺。然后乘之。所以养安也。孰知夫出死要节之所以养生也。孰知夫轻费用之所以养财也。孰知夫恭敬辞让之所以养安也。孰知夫礼义文理之所以养情也。\\
\hspace*{2em}人苟生之为见。若者必死。苟利之为见。若者必害。怠惰之为安。若者必危。情胜之为安。若者必灭。故圣人一之于礼义。则两得之矣。一之于情性。则两失之矣。故儒者将使人两得之者也。墨者将使人两失之者也。是儒墨之分。\\
\hspace*{2em}治辨之极也。强固之本也。威行之道也。功名之总也。王公由之。所以一天下。臣诸侯也。弗由之。所以捐社稷也。故坚革利兵不足以为胜。高城深池不足以为固。严令繁刑不足以为威。由其道则行。不由其道则废。楚人鲛革犀兕。所以为甲。坚如金石。宛之钜铁施。钻如蜂虿。轻利剽遬。卒如熛风。然而兵殆于垂涉。唐昧死焉。庄蹻起。楚分而为四参。是岂无坚革利兵哉。其所以统之者非其道故也。汝颍以为险。江汉以为池。阻之以邓林。缘之以方城。然而秦师至鄢郢。举若振藁。是岂无固塞险阻哉。其所以统之者非其道故也。纣剖比干。囚箕子。为炮格。刑杀无辜。时臣下懔然。莫必其命。然而周师至。而令不行乎下。不能用其民。是岂令不严。刑不陖哉。其所以统之者非其道故也。\\
\hspace*{2em}古者之兵。戈矛弓矢而已。然而敌国不待试而诎。城郭不集。沟池不掘。固塞不树。机变不张。然而国晏然不畏外而固者。无他故焉。明道而均分之。时使而诚爱之。则下应之如景响。有不由命者。然后俟之以刑。则民知罪矣。故刑一人而天下服。罪人不尤其上。知罪之在己也。是故刑罚省而威行如流。无他故焉。由其道故也。故由其道则行。不由其道则废。古者帝尧之治天下也。盖杀一人刑二人而天下治。传曰。﹁威厉而不试。刑措而不用。﹂。\\
\hspace*{2em}天地者。生之本也。先祖者。类之本也。君师者。治之本也。无天地恶生。无先祖恶出。无君师恶治。三者偏亡。则无安人。故礼。上事天。下事地。尊先祖而隆君师。是礼之三本也。\\
\hspace*{2em}故王者天太祖。诸侯不敢怀。大夫士有常宗。所以辨贵贱。贵贱治。得之本也。郊畴乎天子。社至乎诸侯。函及士大夫。所以辨尊者事尊。卑者事卑。宜钜者钜。宜小者小。故有天下者事七世。有一国者事五世。有五乘之地者事三世。有三乘之地者事二世。有特牲而食者不得立宗庙。所以辨积厚者流泽广。积薄者流泽狭也。\\
\hspace*{2em}大飨上玄尊。俎上腥鱼。先大羹。贵食饮之本也。大飨上玄尊而用薄酒。食先黍稷而饭稻粱。祭哜先大羹而饱庶羞。贵本而亲用也。贵本之谓文。亲用之谓理。两者合而成文。以归太一。是谓大隆。故尊之上玄尊也。俎之上腥鱼也。豆之先大羹。一也。利爵弗啐也。成事俎弗尝也。三侑之弗食也。大昏之未废齐也。大庙之未内尸也。始绝之未小敛。一也。大路之素帱也。郊之麻絻。丧服之先散麻。一也。三年哭之不反也。清庙之歌一倡而三叹。县一钟尚拊膈。朱弦而通越。一也。\\
\hspace*{2em}凡礼始乎脱。成乎文。终乎税。故至备。情文俱尽。其次。情文代胜。其下。复情以归太一。天地以合。日月以明。四时以序。星辰以行。江河以流。万物以昌。好恶以节。喜怒以当。以为下则顺。以为上则明。\\
\hspace*{2em}太史公曰。至矣哉。立隆以为极。而天下莫之能益损也。本末相顺。终始相应。至文有以辨。至察有以说。天下从之者治。不从者乱。从之者安。不从者危。小人不能则也。\\
\hspace*{2em}礼之貌诚深矣。坚白同异之察。入焉而弱。其貌诚大矣。擅作典制褊陋之说。入焉而望。其貌诚高矣。暴慢恣睢。轻俗以为高之属。入焉而队。故绳诚陈。则不可欺以曲直。衡诚悬。则不可欺以轻重。规矩诚错。则不可欺以方员。君子审礼。则不可欺以诈伪。故绳者。直之至也。衡者。平之至也。规矩者。方员之至也。礼者。人道之极也。然而不法礼者不足礼。谓之无方之民。法礼足礼。谓之有方之士。礼之中。能思索。谓之能虑。能虑勿易。谓之能固。能虑能固。加好之焉。圣矣。天者。高之极也。地者。下之极也。日月者。明之极也。无穷者。广大之极也。圣人者。道之极也。\\
\hspace*{2em}以财物为用。以贵贱为文。以多少为异。以隆杀为要。文貌繁。情欲省。礼之隆也。文貌省。情欲繁。礼之杀也。文貌情欲相为内外表里。并行而杂。礼之中流也。君子上致其隆。下尽其杀。而中处其中。步骤驰骋广骛不外。是以君子之性守宫庭也。人域是域。士君子也。外是。民也。于是中焉。房皇周浃。曲得其次序。圣人也。故厚者。礼之积也。大者。礼之广也。高者。礼之隆也。明者。礼之尽也。\\

% ===== 索隐述赞（司马贞《史记索隐》卷二十九）=====
\clearpage
\shuzanhead{索隐述赞曰}\linebreak
{\shuzanbody\fontsize{24}{41}\selectfont
\hbox to \hsize{\hfil 礼因人心。非从天下。合诚饰貌。救弊兴雅。\hfil}\linebreak
\hbox to \hsize{\hfil 以制黎甿。以事宗社。情文可重。丰杀难假。\hfil}\linebreak
\hbox to \hsize{\hfil 仲尼坐树。孙通蕝野。圣人作教。罔不由者。\hfil}%
\par}

% ===== 卷二十四　乐书　书口/扉页篇名=乐书 =====
\clearpage
\ExplSyntaxOn
\tl_set:Nn \l__zhvt_chapter_title_tl { 乐书 }
\tl_set:Nn \l__zhvt_title_display_tl { 乐书 }
\ExplSyntaxOff
\maketitle

% 十表八书补排：单篇成册，不渲染目录页




% 方案①：仅改 A 类错讹，头部与版式不变

% ===== 正文（八书续排：vRain 白文底本 024，W 盘中华本核对，造字依定稿口径回改）=====
\juan{乐书}{卷二十四~乐书第二}
\hspace*{2em}太史公曰。余每读虞书。至于君臣相敕。维是几安。而股肱不良。万事堕坏。未尝不流涕也。成王作颂。推己惩艾。悲彼家难。可不谓战战恐惧。善守善终哉。君子不为约则修德。满则弃礼。佚能思初。安能惟始。沐浴膏泽而歌咏勤苦。非大德谁能如斯。传曰﹁治定功成。礼乐乃兴﹂。海内人道益深。其德益至。所乐者益异。满而不损则溢。盈而不持则倾。凡作乐者。所以节乐。君子以谦退为礼。以损减为乐。乐其如此也。以为州异国殊。情习不同。故博采风俗。协比声律。以补短移化。助流政教。天子躬于明堂临观。而万民咸荡涤邪秽。斟酌饱满。以饰厥性。故云雅颂之音理而民正。嘄噭之声兴而士奋。郑卫之曲动而心淫。及其调和谐合。鸟兽尽感。而况怀五常。含好恶。自然之势也。\\
\hspace*{2em}治道亏缺而郑音兴起。封君世辟。名显邻州。争以相高。自仲尼不能与齐优遂容于鲁。虽退正乐以诱世。作五章以刺时。犹莫之化。陵迟以至六国。流沔沈佚。遂往不返。卒于丧身灭宗。并国于秦。\\
\hspace*{2em}秦二世尤以为娱。丞相李斯进谏曰。﹁放弃诗书。极意声色。祖伊所以惧也。轻积细过。恣心长夜。纣所以亡也。﹂赵高曰。﹁五帝。三王乐各殊名。示不相袭。上自朝廷。下至人民。得以接欢喜。合殷勤。非此和说不通。解泽不流。亦各一世之化。度时之乐。何必华山之騄耳而后行远乎。﹂二世然之。\\
\hspace*{2em}高祖过沛诗三侯之章。令小儿歌之。高祖崩。令沛得以四时歌舞宗庙。孝惠。孝文。孝景无所增更。于乐府习常肄旧而已。\\
\hspace*{2em}至今上即位。作十九章。令侍中李延年次序其声。拜为协律都尉。通一经之士不能独知其辞。皆集会五经家。相与共讲习读之。乃能通知其意。多尔雅之文。\\
\hspace*{2em}汉家常以正月上辛祠太一甘泉。以昏时夜祠。到明而终。常有流星经于祠坛上。使僮男僮女七十人俱歌。春歌青阳。夏歌朱明。秋歌西暤。冬歌玄冥。世多有。故不论。\\
\hspace*{2em}又尝得神马渥洼水中。复次以为太一之歌。歌曲曰。﹁太一贡兮天马下。沾赤汗兮沫流赭。骋容与兮跇万里。今安匹兮龙为友。﹂后伐大宛得千里马。马名﹁蒲梢﹂。次作以为歌。歌诗曰。﹁天马来兮从西极。经万里兮归有德。承灵威兮降外国。涉流沙兮四夷服。﹂中尉汲黯进曰。﹁凡王者作乐。上以承祖宗。下以化兆民。今陛下得马。诗以为歌。协于宗庙。先帝百姓岂能知其音邪。﹂上默然不说。丞相公孙弘曰。﹁黯诽谤圣制。当族。﹂。\\
\hspace*{2em}凡音之起。由人心生也。人心之动。物使之然也。感于物而动。故形于声。声相应。故生变。变成方。谓之音。比音而乐之。及干戚羽旄。谓之乐也。乐者。音之所由生也。其本在人心感于物也。是故其哀心感者。其声噍以杀。其乐心感者。其声啴以缓。其喜心感者。其声发以散。其怒心感者。其声粗以厉。其敬心感者。其声直以廉。其爱心感者。其声和以柔。六者非性也。感于物而后动。是故先王慎所以感之。故礼以导其志。乐以和其声。政以壹其行。刑以防其奸。礼乐刑政。其极一也。所以同民心而出治道也。\\
\hspace*{2em}凡音者。生人心者也。情动于中。故形于声。声成文谓之音。是故治世之音安以乐。其正和。乱世之音怨以怒。其正乖。亡国之音哀以思。其民困。声音之道。与正通矣。宫为君。商为臣。角为民。征为事。羽为物。五者不乱。则无惉懘之音矣。宫乱则荒。其君骄。商乱则捶。其臣坏。角乱则忧。其民怨。征乱则哀。其事勤。羽乱则危。其财匮。五者皆乱。迭相陵。谓之慢。如此则国之灭亡无日矣。郑卫之音。乱世之音也。比于慢矣。桑间濮上之音。亡国之音也。其政散。其民流。诬上行私而不可止。\\
\hspace*{2em}凡音者。生于人心者也。乐者。通于伦理者也。是故知声而不知音者。禽兽是也。知音而不知乐者。众庶是也。唯君子为能知乐。是故审声以知音。审音以知乐。审乐以知政。而治道备矣。是故不知声者不可与言音。不知音者不可与言乐。知乐则几于礼矣。礼乐皆得。谓之有德。德者得也。是故乐之隆。非极音也。食飨之礼。非极味也。清庙之瑟。朱弦而疏越。一倡而三叹。有遗音者矣。大飨之礼。尚玄酒而俎腥鱼。大羹不和。有遗味者矣。是故先王之制礼乐也。非以极口腹耳目之欲也。将以教民平好恶而反人道之正也。\\
\hspace*{2em}人生而静。天之性也。感于物而动。性之颂也。物至知知。然后好恶形焉。好恶无节于内。知诱于外。不能反己。天理灭矣。夫物之感人无穷。而人之好恶无节。则是物至而人化物也。人化物也者。灭天理而穷人欲者也。于是有悖逆诈伪之心。有淫佚作乱之事。是故强者胁弱。众者暴寡。知者诈愚。勇者苦怯。疾病不养。老幼孤寡不得其所。此大乱之道也。是故先王制礼乐。人为之节。衰麻哭泣。所以节丧纪也。钟鼓干戚。所以和安乐也。婚姻冠笄。所以别男女也。射乡食飨。所以正交接也。礼节民心。乐和民声。政以行之。刑以防之。礼乐刑政四达而不悖。则王道备矣。\\
\hspace*{2em}乐者为同。礼者为异。同则相亲。异则相敬。乐胜则流。礼胜则离。合情饰貌者。礼乐之事也。礼义立。则贵贱等矣。乐文同。则上下和矣。好恶着。则贤不肖别矣。刑禁暴。爵举贤。则政均矣。仁以爱之。义以正之。如此则民治行矣。\\
\hspace*{2em}乐由中出。礼自外作。乐由中出。故静。礼自外作。故文。大乐必易。大礼必简。乐至则无怨。礼至则不争。揖让而治天下者。礼乐之谓也。暴民不作。诸侯宾服。兵革不试。五刑不用。百姓无患。天子不怒。如此则乐达矣。合父子之亲。明长幼之序。以敬四海之内。天子如此。则礼行矣。\\
\hspace*{2em}大乐与天地同和。大礼与天地同节。和。故百物不失。节。故祀天祭地。明则有礼乐。幽则有鬼神。如此则四海之内合敬同爱矣。礼者。殊事合敬者也。乐者。异文合爱者也。礼乐之情同。故明王以相沿也。故事与时并。名与功偕。故钟鼓管磬羽龠干戚。乐之器也。诎信俯仰级兆舒疾。乐之文也。簠簋俎豆制度文章。礼之器也。升降上下周旋裼袭。礼之文也。故知礼乐之情者能作。识礼乐之文者能术。作者之谓圣。术者之谓明。明圣者。术作之谓也。\\
\hspace*{2em}乐者。天地之和也。礼者。天地之序也。和。故百物皆化。序。故群物皆别。乐由天作。礼以地制。过制则乱。过作则暴。明于天地。然后能兴礼乐也。论伦无患。乐之情也。欣喜欢爱。乐之官也。中正无邪。礼之质也。庄敬恭顺。礼之制也。若夫礼乐之施于金石。越于声音。用于宗庙社稷。事于山川鬼神。则此所以与民同也。\\
\hspace*{2em}王者功成作乐。治定制礼。其功大者其乐备。其治辨者其礼具。干戚之舞。非备乐也。亨孰而祀。非达礼也。五帝殊时。不相沿乐。三王异世。不相袭礼。乐极则忧。礼粗则偏矣。及夫敦乐而无忧。礼备而不偏者。其唯大圣乎。天高地下。万物散殊。而礼制行也。流而不息。合同而化。而乐兴也。春作夏长。仁也。秋敛冬藏。义也。仁近于乐。义近于礼。乐者敦和。率神而从天。礼者辨宜。居鬼而从地。故圣人作乐以应天。作礼以配地。礼乐明备。天地官矣。\\
\hspace*{2em}天尊地卑。君臣定矣。高卑已陈。贵贱位矣。动静有常。大小大殊矣。方以类聚。物以群分。则性命不同矣。在天成象。在地成形。如此则体者天地之别也。地气上隮。天气下降。阴阳相摩。天地相荡。鼓之以雷霆。奋之以风雨。动之以四时。暖之以日月。而百化兴焉。如此则乐者天地之和也。\\
\hspace*{2em}化不时则不生。男女无别则乱登。此天地之情也。及夫礼乐之极乎天而蟠乎地。行乎阴阳而通乎鬼神。穷高极远而测深厚。乐着太始而礼居成物。着不息者天也。着不动者地也。一动一静者。天地之间也。故圣人曰﹁礼云乐云﹂。\\
\hspace*{2em}昔者舜作五弦之琴。以歌南风。夔始作乐。以赏诸侯。故天子之为乐也。以赏诸侯之有德者也。德盛而教尊。五谷时孰。然后赏之以乐。故其治民劳者。其舞行级远。其治民佚者。其舞行级短。故观其舞而知其德。闻其谥而知其行。大章。章之也。咸池。备也。韶。继也。夏。大也。殷。周之乐尽也。\\
\hspace*{2em}天地之道。寒暑不时则疾。风雨不节则饥。教者。民之寒暑也。教不时则伤世。事者。民之风雨也。事不节则无功。然则先王之为乐也。以法治也。善则行象德矣。夫豢豕为酒。非以为祸也。而狱讼益烦。则酒之流生祸也。是故先王因为酒礼。一献之礼。宾主百拜。终日饮酒而不得醉焉。此先王之所以备酒祸也。故酒食者。所以合欢也。\\
\hspace*{2em}乐者。所以象德也。礼者。所以闭淫也。是故先王有大事。必有礼以哀之。有大福。必有礼以乐之。哀乐之分。皆以礼终。\\
\hspace*{2em}乐也者。施也。礼也者。报也。乐。乐其所自生。而礼。反其所自始。乐章德。礼报情反始也。所谓大路者。天子之舆也。龙旗九旒。天子之旌也。青黑缘者。天子之葆龟也。从之以牛羊之群。则所以赠诸侯也。\\
\hspace*{2em}乐也者。情之不可变者也。礼也者。理之不可易者也。乐统同。礼别异。礼乐之说贯乎人情矣。穷本知变。乐之情也。着诚去伪。礼之经也。礼乐顺天地之诚。达神明之德。降兴上下之神。而凝是精粗之体。领父子君臣之节。\\
\hspace*{2em}是故大人举礼乐。则天地将为昭焉。天地欣合。阴阳相得。煦妪复育万物。然后草木茂。区萌达。羽翮奋。角觡生。蛰虫昭稣。羽者妪伏。毛者孕鬻。胎生者不殰而卵生者不殈。则乐之道归焉耳。\\
\hspace*{2em}乐者。非谓黄钟大吕弦歌干扬也。乐之末节也。故童者舞之。布筵席。陈樽俎。列笾豆。以升降为礼者。礼之末节也。故有司掌之。乐师辩乎声诗。故北面而弦。宗祝辩乎宗庙之礼。故后尸。商祝辩乎丧礼。故后主人。是故德成而上。蓺成而下。行成而先。事成而后。是故先王有上有下。有先有后。然后可以有制于天下也。\\
\hspace*{2em}乐者。圣人之所乐也。而可以善民心。其感人深。其风移俗易。故先王着其教焉。\\
\hspace*{2em}夫人有血气心知之性。而无哀乐喜怒之常。应感起物而动。然后心术形焉。是故志微焦衰之音作。而民思忧。啴缓慢易繁文简节之音作。而民康乐。粗厉勐起奋末广贲之音作。而民刚毅。廉直经正庄诚之音作。而民肃敬。宽裕肉好顺成和动之音作。而民慈爱。流辟邪散狄成涤滥之音作。而民淫乱。\\
\hspace*{2em}是故先王本之情性。稽之度数。制之礼义。合生气之和。道五常之行。使之阳而不散。阴而不密。刚气不怒。柔气不慑。四畅交于中而发作于外。皆安其位而不相夺也。然后立之学等。广其节奏。省其文采。以绳德厚也。类小大之称。比终始之序。以象事行。使亲疏贵贱长幼男女之理皆形见于乐。故曰﹁乐观其深矣﹂。\\
\hspace*{2em}土敝则草木不长。水烦则鱼鳖不大。气衰则生物不育。世乱则礼废而乐淫。是故其声哀而不庄。乐而不安。慢易以犯节。流湎以忘本。广则容奸。狭则思欲。感涤荡之气而灭平和之德。是以君子贱之也。\\
\hspace*{2em}凡奸声感人而逆气应之。逆气成象而淫乐兴焉。正声感人而顺气应之。顺气成象而和乐兴焉。倡和有应。回邪曲直各归其分。而万物之理以类相动也。\\
\hspace*{2em}是故君子反情以和其志。比类以成其行。奸声乱色不留聪明。淫乐废礼不接于心术。惰慢邪辟之气不设于身体。使耳目鼻口心知百体皆由顺正。以行其义。然后发以声音。文以琴瑟。动以干戚。饰以羽旄。从以箫管。奋至德之光。动四气之和。以着万物之理。是故清明象天。广大象地。终始象四时。周旋象风雨。五色成文而不乱。八风从律而不奸。百度得数而有常。小大相成。终始相生。倡和清浊。代相为经。故乐行而伦清。耳目聪明。血气和平。移风易俗。天下皆宁。故曰﹁乐者乐也﹂。君子乐得其道。小人乐得其欲。以道制欲。则乐而不乱。以欲忘道。则惑而不乐。是故君子反情以和其志。广乐以成其教。乐行而民乡方。可以观德矣。\\
\hspace*{2em}德者。性之端也。乐者。德之华也。金石丝竹。乐之器也。诗。言其志也。歌。咏其声也。舞。动其容也。三者本乎心。然后乐气从之。是故情深而文明。气盛而化神。和顺积中而英华发外。唯乐不可以为伪。\\
\hspace*{2em}乐者。心之动也。声者。乐之象也。文采节奏。声之饰也。君子动其本。乐其象。然后治其饰。是故先鼓以警戒。三步以见方。再始以着往。复乱以饬归。奋疾而不拔。极幽而不隐。独乐其志。不厌其道。备举其道。不私其欲。是以情见而义立。乐终而德尊。君子以好善。小人以息过。故曰﹁生民之道。乐为大焉﹂。\\
\hspace*{2em}君子曰。礼乐不可以斯须去身。致乐以治心。则易直子谅之心油然生矣。易直子谅之心生则乐。乐则安。安则久。久则天。天则神。天则不言而信。神则不怒而威。致乐。以治心者也。致礼。以治躬者也。治躬则庄敬。庄敬则严威。心中斯须不和不乐。而鄙诈之心入之矣。外貌斯须不庄不敬。而慢易之心入之矣。故乐也者。动于内者也。礼也者。动于外者也。乐极和。礼极顺。内和而外顺。则民瞻其颜色而弗与争也。望其容貌而民不生易慢焉。德辉动乎内而民莫不承听。理发乎外而民莫不承顺。故曰﹁知礼乐之道。举而错之天下无难矣﹂。\\
\hspace*{2em}乐也者。动于内者也。礼也者。动于外者也。故礼主其谦。乐主其盈。礼谦而进。以进为文。乐盈而反。以反为文。礼谦而不进。则销。乐盈而不反。则放。故礼有报而乐有反。礼得其报则乐。乐得其反则安。礼之报。乐之反。其义一也。\\
\hspace*{2em}夫乐者乐也。人情之所不能免也。乐必发诸声音。形于动静。人道也。声音动静。性术之变。尽于此矣。故人不能无乐。乐不能无形。形而不为道。不能无乱。先王恶其乱。故制雅颂之声以道之。使其声足以乐而不流。使其文足以纶而不息。使其曲直繁省廉肉节奏。足以感动人之善心而已矣。不使放心邪气得接焉。是先王立乐之方也。是故乐在宗庙之中。君臣上下同听之。则莫不和敬。在族长乡里之中。长幼同听之。则莫不和顺。在闺门之内。父子兄弟同听之。则莫不和亲。故乐者。审一以定和。比物以饰节。节奏合以成文。所以合和父子君臣。附亲万民也。是先王立乐之方也。故听其雅颂之声。志意得广焉。执其干戚。习其俯仰诎信。容貌得庄焉。行其缀兆。要其节奏。行列得正焉。进退得齐焉。故乐者天地之齐。中和之纪。人情之所不能免也。\\
\hspace*{2em}夫乐者。先王之所以饰喜也。军旅鈇钺者。先王之所以饰怒也。故先王之喜怒皆得其齐矣。喜则天下和之。怒则暴乱者畏之。先王之道礼乐可谓盛矣。\\
\hspace*{2em}魏文侯问于子夏曰。﹁吾端冕而听古乐则唯恐卧。听郑卫之音则不知倦。敢问古乐之如彼。何也。新乐之如此。何也。﹂。\\
\hspace*{2em}子夏答曰。﹁今夫古乐。进旅而退旅。和正以广。弦匏笙簧合守拊鼓。始奏以文。止乱以武。治乱以相。讯疾以雅。君子于是语。于是道古。修身及家。平均天下。此古乐之发也。今夫新乐。进俯退俯。奸声以淫。溺而不止。及优侏儒。獶杂子女。不知父子。乐终不可以语。不可以道古。此新乐之发也。今君之所问者乐也。所好者音也。夫乐之与音。相近而不同。﹂。\\
\hspace*{2em}文侯曰。﹁敢问如何。﹂。\\
\hspace*{2em}子夏答曰。﹁夫古者天地顺而四时当。民有德而五谷昌。疾疢不作而无祆祥。此之谓大当。然后圣人作为父子君臣以为之纪纲。纪纲既正。天下大定。天下大定。然后正六律。和五声。弦歌诗颂。此之谓德音。德音之谓乐。诗曰。﹁莫其德音。其德克明。克明克类。克长克君。王此大邦。克顺克俾。俾于文王。其德靡悔。既受帝祉。施于孙子。﹂此之谓也。今君之所好者。其溺音与。﹂。\\
\hspace*{2em}文侯曰。﹁敢问溺音者何从出也。﹂。\\
\hspace*{2em}子夏答曰。﹁郑音好滥淫志。宋音燕女溺志。卫音趣数烦志。齐音骜辟骄志。四者皆淫于色而害于德。是以祭祀不用也。诗曰。﹁肃雍和鸣。先祖是听。﹂夫肃肃。敬也。雍雍。和也。夫敬以和。何事不行。为人君者。谨其所好恶而已矣。君好之则臣为之。上行之则民从之。诗曰。﹁诱民孔易﹂。此之谓也。然后圣人作为鼗鼓椌楬埙篪。此六者。德音之音也。然后钟磬竽瑟以和之。干戚旄狄以舞之。此所以祭先王之庙也。所以献酬酳酢也。所以官序贵贱各得其宜也。此所以示后世有尊卑长幼序也。钟声铿。铿以立号。号以立横。横以立武。君子听钟声则思武臣。石声硁。硁以立别。别以致死。君子听磬声则思死封疆之臣。丝声哀。哀以立廉。廉以立志。君子听琴瑟之声则思志义之臣。竹声滥。滥以立会。会以聚众。君子听竽笙箫管之声则思畜聚之臣。鼓鼙之声欢。欢以立动。动以进众。君子听鼓鼙之声则思将帅之臣。君子之听音。非听其铿锵而已也。彼亦有所合之也。﹂。\\
\hspace*{2em}宾牟贾侍坐于孔子。孔子与之言。及乐。曰。﹁夫武之备戒之已久。何也。﹂。\\
\hspace*{2em}答曰。﹁病不得其众也。﹂。\\
\hspace*{2em}﹁永叹之。淫液之。何也。﹂。\\
\hspace*{2em}答曰。﹁恐不逮事也。﹂。\\
\hspace*{2em}﹁发扬蹈厉之已蚤。何也。﹂。\\
\hspace*{2em}答曰。﹁及时事也。﹂。\\
\hspace*{2em}﹁武坐致右宪左。何也。﹂。\\
\hspace*{2em}答曰。﹁非武坐也。﹂。\\
\hspace*{2em}﹁声淫及商。何也。﹂。\\
\hspace*{2em}答曰。﹁非武音也。﹂。\\
\hspace*{2em}子曰。﹁若非武音。则何音也。﹂。\\
\hspace*{2em}答曰。﹁有司失其传也。如非有司失其传。则武王之志荒矣。﹂。\\
\hspace*{2em}子曰。﹁唯丘之闻诸苌弘。亦若吾子之言是也。﹂。\\
\hspace*{2em}宾牟贾起。免席而请曰。﹁夫武之备戒之已久。则既闻命矣。敢问迟之迟而又久。何也。﹂。\\
\hspace*{2em}子曰。﹁居。吾语汝。夫乐者。象成者也。总干而山立。武王之事也。发扬蹈厉。太公之志也。武乱皆坐。周召之治也。且夫武。始而北出。再成而灭商。三成而南。四成而南国是疆。五成而分陕。周公左。召公右。六成复缀。以崇天子。夹振之而四伐。盛威于中国也。分夹而进。事蚤济也。久立于缀。以待诸侯之至也。且夫女独未闻牧野之语乎。武王克殷反商。未及下车。而封黄帝之后于蓟。封帝尧之后于祝。封帝舜之后于陈。下车而封夏后氏之后于杞。封殷之后于宋。封王子比干之墓。释箕子之囚。使之行商容而复其位。庶民弛政。庶士倍禄。济河而西。马散华山之阳而弗复乘。牛散桃林之野而不复服。车甲弢而藏之府库而弗复用。倒载干戈。苞之以虎皮。将率之士。使为诸侯。名之曰﹁建櫜﹂。然后天下知武王之不复用兵也。散军而郊射。左射狸首。右射驺虞。而贯革之射息也。裨冕搢笏。而虎贲之士税剑也。祀乎明堂。而民知孝。朝觐。然后诸侯知所以臣。耕藉。然后诸侯知所以敬。五者天下之大教也。食三老五更于太学。天子袒而割牲。执酱而馈。执爵而酳。冕而总干。所以教诸侯之悌也。若此。则周道四达。礼乐交通。则夫武之迟久。不亦宜乎。﹂。\\
\hspace*{2em}子贡见师乙而问焉。曰。﹁赐闻声歌各有宜也。如赐者宜何歌也。﹂。\\
\hspace*{2em}师乙曰。﹁乙。贱工也。何足以问所宜。请诵其所闻。而吾子自执焉。宽而静。柔而正者宜歌颂。广大而静。疏达而信者宜歌大雅。恭俭而好礼者宜歌小雅。正直清廉而谦者宜歌风。肆直而慈爱者宜歌商。温良而能断者宜歌齐。夫歌者。直己而陈德。动己而天地应焉。四时和焉。星辰理焉。万物育焉。故商者。五帝之遗声也。商人志之。故谓之商。齐者。三代之遗声也。齐人志之。故谓之齐。明乎商之诗者。临事而屡断。明乎齐之诗者。见利而让也。临事而屡断。勇也。见利而让。义也。有勇有义。非歌孰能保此。故歌者。上如抗。下如队。曲如折。止如藁木。居中矩。句中钩。累累乎殷如贯珠。故歌之为言也。长言之也。说之。故言之。言之不足。故长言之。长言之不足。故嗟叹之。嗟叹之不足。故不知手之舞之足之蹈之。﹂子贡问乐。\\
\hspace*{2em}凡音由于人心。天之与人有以相通。如景之象形。响之应声。故为善者天报之以福。为恶者天与之以殃。其自然者也。\\
\hspace*{2em}故舜弹五弦之琴。歌南风之诗而天下治。纣为朝歌北鄙之音。身死国亡。舜之道何弘也。纣之道何隘也。夫南风之诗者生长之音也。舜乐好之。乐与天地同意。得万国之心。故天下治也。夫朝歌者不时也。北者败也。鄙者陋也。纣乐好之。与万国殊心。诸侯不附。百姓不亲。天下畔之。故身死国亡。\\
\hspace*{2em}而卫灵公之时。将之晋。至于濮水之上舍。夜半时闻鼓琴声。问左右。皆对曰﹁不闻﹂。乃召师涓曰。﹁吾闻鼓琴音。问左右。皆不闻。其状似鬼神。为我听而写之。﹂师涓曰。﹁诺。﹂因端坐援琴。听而写之。明日。曰。﹁臣得之矣。然未习也。请宿习之。﹂灵公曰。﹁可。﹂因复宿。明日。报曰。﹁习矣。﹂即去之晋。见晋平公。平公置酒于施惠之台。酒酣。灵公曰。﹁今者来。闻新声。请奏之。﹂平公曰。﹁可。﹂即令师涓坐师旷旁。援琴鼓之。未终。师旷抚而止之曰。﹁此亡国之声也。不可遂。﹂平公曰。﹁何道出。﹂师旷曰。﹁师延所作也。与纣为靡靡之乐。武王伐纣。师延东走。自投濮水之中。故闻此声必于濮水之上。先闻此声者国削。﹂平公曰。﹁寡人所好者音也。愿遂闻之。﹂师涓鼓而终之。\\
\hspace*{2em}平公曰。﹁音无此最悲乎。﹂师旷曰。﹁有。﹂平公曰。﹁可得闻乎。﹂师旷曰。﹁君德义薄。不可以听之。﹂平公曰。﹁寡人所好者音也。愿闻之。﹂师旷不得已。援琴而鼓之。一奏之。有玄鹤二八集乎廊门。再奏之。延颈而鸣。舒翼而舞。\\
\hspace*{2em}平公大喜。起而为师旷寿。反坐。问曰。﹁音无此最悲乎。﹂师旷曰。﹁有。昔者黄帝以大合鬼神。今君德义薄。不足以听之。听之将败。﹂平公曰。﹁寡人老矣。所好者音也。愿遂闻之。﹂师旷不得已。援琴而鼓之。一奏之。有白云从西北起。再奏之。大风至而雨随之。飞廊瓦。左右皆奔走。平公恐惧。伏于廊屋之间。晋国大旱。赤地三年。\\
\hspace*{2em}听者或吉或凶。夫乐不可妄兴也。\\
\hspace*{2em}太史公曰。夫上古明王举乐者。非以娱心自乐。快意恣欲。将欲为治也。正教者皆始于音。音正而行正。故音乐者。所以动荡血脉。通流精神而和正心也。故宫动脾而和正圣。商动肺而和正义。角动肝而和正仁。徵动心而和正礼。羽动肾而和正智。故乐所以内辅正心而外异贵贱也。上以事宗庙。下以变化黎庶也。琴长八尺一寸。正度也。弦大者为宫。而居中央。君也。商张右傍。其余大小相次。不失其次序。则君臣之位正矣。故闻宫音。使人温舒而广大。闻商音。使人方正而好义。闻角音。使人恻隐而爱人。闻徵音。使人乐善而好施。闻羽音。使人整齐而好礼。夫礼由外入。乐自内出。故君子不可须臾离礼。须臾离礼则暴慢之行穷外。不可须臾离乐。须臾离乐则奸邪之行穷内。故乐音者。君子之所养义也。夫古者。天子诸侯听钟磬未尝离于庭。卿大夫听琴瑟之音未尝离于前。所以养行义而防淫佚也。夫淫佚生于无礼。故圣王使人耳闻雅颂之音。目视威仪之礼。足行恭敬之容。口言仁义之道。故君子终日言而邪辟无由入也。\\

% ===== 索隐述赞（司马贞《史记索隐》卷二十九）=====
\clearpage
\shuzanhead{索隐述赞曰}\linebreak
{\shuzanbody\fontsize{24}{41}\selectfont
\hbox to \hsize{\hfil 乐之所兴。在乎防欲。陶心畅志。舞手蹈足。\hfil}\linebreak
\hbox to \hsize{\hfil 舜曰箫韶。融称属续。审音知政。观风变俗。\hfil}\linebreak
\hbox to \hsize{\hfil 端如贯珠。清同叩玉。洋洋盈耳。咸英余曲。\hfil}%
\par}

% ===== 卷二十五　律书　书口/扉页篇名=律书 =====
\clearpage
\ExplSyntaxOn
\tl_set:Nn \l__zhvt_chapter_title_tl { 律书 }
\tl_set:Nn \l__zhvt_title_display_tl { 律书 }
\ExplSyntaxOff
\maketitle

% 十表八书补排：单篇成册，不渲染目录页




% 方案①：仅改 A 类错讹，头部与版式不变

% ===== 正文（八书续排：vRain 白文底本 025，W 盘中华本核对，造字依定稿口径回改）=====
\juan{律书}{卷二十五~律书第三}
\hspace*{2em}王者制事立法。物度轨则。壹禀于六律。六律为万事根本焉。\\
\hspace*{2em}其于兵械尤所重。故云﹁望敌知吉凶。闻声效胜负﹂。百王不易之道也。\\
\hspace*{2em}武王伐纣。吹律听声。推孟春以至于季冬。杀气相并。而音尚宫。同声相从。物之自然。何足怪哉。\\
\hspace*{2em}兵者。圣人所以讨强暴。平乱世。夷险阻。救危殆。自含齿戴角之兽见犯则校。而况于人怀好恶喜怒之气。喜则爱心生。怒则毒螫加。情性之理也。\\
\hspace*{2em}昔黄帝有涿鹿之战。以定火灾。颛顼有共工之陈。以平水害。成汤有南巢之伐。以殄夏乱。递兴递废。胜者用事。所受于天也。\\
\hspace*{2em}自是之后。名士迭兴。晋用咎犯。而齐用王子。吴用孙武。申明军约。赏罚必信。卒伯诸侯。兼列邦土。虽不及三代之诰誓。然身宠君尊。当世显扬。可不谓荣焉。岂与世儒暗于大较。不权轻重。猥云德化。不当用兵。大至君辱失守。小乃侵犯削弱。遂执不移等哉。故教笞不可废于家。刑罚不可捐于国。诛伐不可偃于天下。用之有巧拙。行之有逆顺耳。\\
\hspace*{2em}夏桀。殷纣手搏豺狼。足追四马。勇非微也。百战克胜。诸侯慑服。权非轻也。秦二世宿军无用之地。连兵于边陲。力非弱也。结怨匈奴。祸于越。势非寡也。及其威尽势极。闾巷之人为敌国。咎生穷武之不知足。甘得之心不息也。\\
\hspace*{2em}高祖有天下。三边外畔。大国之王虽称蕃辅。臣节未尽。会高祖厌苦军事。亦有萧。张之谋。故偃武一休息。羁縻不备。\\
\hspace*{2em}历至孝文即位。将军陈武等议曰。﹁南越。朝鲜自全秦时内属为臣子。后且拥兵阻阸。选蠕观望。高祖时天下新定。人民小安。未可复兴兵。今陛下仁惠抚百姓。恩泽加海内。宜及士民乐用。征讨逆党。以一封疆。﹂孝文曰。﹁朕能任衣冠。念不到此。会吕氏之乱。功臣宗室共不羞耻。误居正位。常战战栗栗。恐事之不终。且兵凶器。虽克所愿。动亦秏病。谓百姓远方何。又先帝知劳民不可烦。故不以为意。朕岂自谓能。今匈奴内侵。军吏无功。边民父子荷兵日久。朕常为动心伤痛。无日忘之。今未能销距。愿且坚边设侯。结和通使。休宁北陲。为功多矣。且无议军。﹂故百姓无内外之繇。得息肩于田亩。天下殷富。粟至十余钱。鸣鸡吠狗。烟火万里。可谓和乐者乎。\\
\hspace*{2em}太史公曰。文帝时。会天下新去汤火。人民乐业。因其欲然。能不扰乱。故百姓遂安。自年六七十翁亦未尝至市井。游敖嬉戏如小儿状。孔子所称有德君子者邪。\\
\hspace*{2em}书曰﹁七正﹂。二十八舍。律历。天所以通五行八正之气。天所以成孰万物也。舍者。日月所舍。舍者。舒气也。\\
\hspace*{2em}不周风居西北。主杀生。东壁居不周风东。主辟生气而东之。至于营室。营室者。主营胎阳气而产之。东至于危。危。垝也。言阳气之垝。故曰危。十月也。律中应钟。应钟者。阳气之应。不用事也。其于十二子为亥。亥者。该也。言阳气藏于下。故该也。\\
\hspace*{2em}广莫风居北方。广莫者。言阳气在下。阴莫阳广大也。故曰广莫。东至于虚。虚者。能实能虚。言阳气冬则宛藏于虚。日冬至则一阴下藏。一阳上舒。故曰虚。东至于须女。言万物变动其所。阴阳气未相离。尚相胥如也。故曰须女。十一月也。律中黄钟。黄钟者。阳气踵黄泉而出也。其于十二子为子。子者。滋也。滋者。言万物滋于下也。其于十母为壬癸。壬之为言任也。言阳气任养万物于下也。癸之为言揆也。言万物可揆度。故曰癸。东至牵牛。牵牛者。言阳气牵引万物出之也。牛者。冒也。言地虽冻。能冒而生也。牛者。耕植种万物也。东至于建星。建星者。建诸生也。十二月也。律中大吕。大吕者。其于十二子为丑。\\
\hspace*{2em}条风居东北。主出万物。条之言条治万物而出之。故曰条风。南至于箕。箕者。言万物根棋。故曰箕。正月也。律中泰蔟。泰蔟者。言万物蔟生也。故曰泰蔟。其于十二子为寅。寅言万物始生蚓然也。故曰寅。南至于尾。言万物始生如尾也。南至于心。言万物始生有华心也。南至于房。房者。言万物门户也。至于门则出矣。\\
\hspace*{2em}明庶风居东方。明庶者。明众物尽出也。二月也。律中夹钟。夹钟者。言阴阳相夹厕也。其于十二子为卯。卯之为言茂也。言万物茂也。其于十母为甲乙。甲者。言万物剖符甲而出也。乙者。言万物生轧轧也。南至于氐。氐者。言万物皆至也。南至于亢。亢者。言万物亢见也。南至于角。角者。言万物皆有枝格如角也。三月也。律中姑洗。姑洗者。言万物洗生。其于十二子为辰。辰者。言万物之蜄也。\\
\hspace*{2em}清明风居东南维。主风吹万物而西之。至于轸。轸者。言万物益大而轸轸然。西至于翼。翼者。言万物皆有羽翼也。四月也。律中中吕。中吕者。言万物尽旅而西行也。其于十二子为巳。巳者。言阳气之已尽也。西至于七星。七星者。阳数成于七。故曰七星。西至于张。张者。言万物皆张也。西至于注。注者。言万物之始衰。阳气下注。故曰注。五月也。律中蕤宾。蕤宾者。言阴气幼少。故曰蕤。痿阳不用事。故曰宾。\\
\hspace*{2em}景风居南方。景者。言阳气道竟。故曰景风。其于十二子为午。午者。阴阳交。故曰午。其于十母为丙丁。丙者。言阳道着明。故曰丙。丁者。言万物之丁壮也。故曰丁。西至于弧。弧者。言万物之吴落且就死也。西至于狼。狼者。言万物可度量。断万物。故曰狼。\\
\hspace*{2em}凉风居西南维。主地。地者。沉夺万物气也。六月也。律中林钟。林钟者。言万物就死气林林然。其于十二子为未。未者。言万物皆成。有滋味也。北至于罚。罚者。言万物气夺可伐也。北至于参。参言万物可参也。故曰参。七月也。律中夷则。夷则。言阴气之贼万物也。其于十二子为申。申者。言阴用事。申贼万物。故曰申。北至于浊。浊者。触也。言万物皆触死也。故曰浊。北至于留。留者。言阳气之稽留也。故曰留。八月也。律中南吕。南吕者。言阳气之旅入藏也。其于十二子为酉。酉者。万物之老也。故曰酉。\\
\hspace*{2em}阊阖风居西方。阊者。倡也。阖者。藏也。言阳气道万物。阖黄泉也。其于十母为庚辛。庚者。言阴气庚万物。故曰庚。辛者。言万物之辛生。故曰辛。北至于胃。胃者。言阳气就藏。皆胃胃也。北至于娄。娄者。呼万物且内之也。北至于奎。奎者。主毒螫杀万物也。奎而藏之。九月也。律中无射。无射者。阴气盛用事。阳气无余也。故曰无射。其于十二子为戌。戌者。言万物尽灭。故曰戌。\\
\hspace*{2em}律数。九九八十一以为宫。三分去一。五十四以为徵。三分益一。七十二以为商。三分去一。四十八以为羽。三分益一。六十四以为角。黄钟长八寸十分一。宫。大吕长七寸五分三分二。太蔟长七寸十分二。角。夹钟长六寸七分三分一。姑洗长六寸十分四。羽。仲吕长五寸九分三分二。征。蕤宾长五寸六分三分二。林钟长五寸十分四。角。夷则长五寸三分二。商。南吕长四寸十分八。征。无射长四寸四分三分二。应钟长四寸二分三分二。羽。\\
\hspace*{2em}生钟分。子一分。丑三分二。寅九分八。卯二十七分十六。辰八十一分六十四。巳二百四十三分一百二十八。午七百二十九分五百一十二。未二千一百八十七分一千二十四。申六千五百六十一分四千九十六。酉一万九千六百八十三分八千一百九十二。戌五万九千四十九分三万二千七百六十八。亥十七万七千一百四十七分六万五千五百三十六。\\
\hspace*{2em}生黄钟术曰。以下生者。倍其实。三其法。以上生者。四其实。三其法。上九。商八。羽七。角六。宫五。征九。置一而九三之以为法。实如法。得长一寸。凡得九寸。命曰﹁黄钟之宫﹂。故曰音始于宫。穷于角。数始于一。终于十。成于三。气始于冬至。周而复生。\\
\hspace*{2em}神生于无。形成于有。形然后数。形而成声。故曰神使气。气就形。形理如类有可类。或未形而未类。或同形而同类。类而可班。类而可识。圣人知天地识之别。故从有以至未有。以得细若气。微若声。然圣人因神而存之。虽妙必效情。核其华道者明矣。非有圣心以乘聪明。孰能存天地之神而成形之情哉。神者。物受之而不能知其去来。故圣人畏而欲存之。唯欲存之。神之亦存。其欲存之者。故莫贵焉。\\
\hspace*{2em}太史公曰。在旋玑玉衡以齐七政。即天地。二十八宿。十母。十二子。钟律调自上古。建律运历造日度。可据而度也。合符节。信道德。即从斯之谓也。\\

% ===== 索隐述赞（司马贞《史记索隐》卷二十九）=====
\clearpage
\shuzanhead{索隐述赞曰}\linebreak
{\shuzanbody\fontsize{24}{41}\selectfont
\hbox to \hsize{\hfil 自昔轩后。爰命伶伦。雄雌是听。厚薄伊均。\hfil}\linebreak
\hbox to \hsize{\hfil 以调气候。以轨星辰。军容取节。乐器斯因。\hfil}\linebreak
\hbox to \hsize{\hfil 自微至著。测化穷神。大哉虚受。含养生人。\hfil}%
\par}

% ===== 卷二十六　历书　书口/扉页篇名=历书 =====
\clearpage
\ExplSyntaxOn
\tl_set:Nn \l__zhvt_chapter_title_tl { 历书 }
\tl_set:Nn \l__zhvt_title_display_tl { 历书 }
\ExplSyntaxOff
\maketitle

% 十表八书补排：单篇成册，不渲染目录页




% 方案①：仅改 A 类错讹，头部与版式不变

% ===== 正文（八书续排：vRain 白文底本 026，W 盘中华本核对，造字依定稿口径回改）=====
\juan{历书}{卷二十六~历书第四}
\hspace*{2em}昔自在古。历建正作於孟春。於时冰泮发蛰。百草奋兴。秭鳺先滜。物乃岁具。生於东。次顺四时。卒于冬分。时鸡三号。卒明。抚十二节。卒于丑。日月成。故明也。明者孟也。幽者幼也。幽明者雌雄也。雌雄代兴。而顺至正之统也。日归于西。起明於东。月归於东。起明于西。正不率天。又不由人。则凡事易坏而难成矣。\\
\hspace*{2em}王者易姓受命。必慎始初。改正朔。易服色。推本天元。顺承厥意。\\
\hspace*{2em}太史公曰。神农以前尚矣。盖黄帝考定星历。建立五行。起消息。正闰馀。於是有天地神祇物类之官。是谓五官。各司其序。不相乱也。民是以能有信。神是以能有明德。民神异业。敬而不渎。故神降之嘉生。民以物享。灾祸不生。所求不匮。\\
\hspace*{2em}少昚氏之衰也。九黎乱德。民神杂扰。不可放物。祸菑荐至。莫尽其气。颛顼受之。乃命南正重司天以属神。命火正黎司地以属民。使复旧常。无相侵渎。其後三苗服九黎之德。故二官咸废所职。而闰馀乖次。孟陬殄灭。摄提无纪。历数失序。尧复遂重黎之後。不忘旧者。使复典之。而立羲和之官。明时正度。则阴阳调。风雨节。茂气至。民无夭疫。年耆禅舜。申戒文祖。云﹁天之历数在尔躬﹂。舜亦以命禹。由是观之。王者所重也。\\
\hspace*{2em}夏正以正月。殷正以十二月。周正以十一月。盖三王之正若循环。穷则反本。天下有道。则不失纪序。无道。则正朔不行於诸侯。\\
\hspace*{2em}幽。厉之後。周室微。陪臣执政。史不记时。君不告朔。故畴人子弟分散。或在诸夏。或在夷狄。是以其禨祥废而不统。周襄王二十六年闰三月。而春秋非之。先王之正时也。履端於始。举正於中。归邪於终。履端於始。序则不愆。举正於中。民则不惑。归邪於终。事则不悖。\\
\hspace*{2em}其後战国并争。在於彊国禽敌。救急解纷而已。岂遑念斯哉。是时独有邹衍。明於五德之传。而散消息之分。以显诸侯。而亦因秦灭六国。兵戎极烦。又升至尊之日浅。未暇遑也。而亦颇推五胜。而自以为获水德之瑞。更名河曰﹁德水﹂。而正以十月。色上黑。然历度闰馀。未能睹其真也。\\
\hspace*{2em}汉兴。高祖曰﹁北畤待我而起﹂。亦自以为获水德之瑞。虽明习历及张苍等。咸以为然。是时天下初定。方纲纪大基。高后女主。皆未遑。故袭秦正朔服色。\\
\hspace*{2em}至孝文时。鲁人公孙臣以终始五德上书。言﹁汉得土德。宜更元。改正朔。易服色。当有瑞。瑞黄龙见﹂。事下丞相张苍。张苍亦学律历。以为非是。罢之。其後黄龙见成纪。张苍自黜。所欲论著不成。而新垣平以望气见。颇言正历服色事。贵幸。後作乱。故孝文帝废不复问。\\
\hspace*{2em}至今上即位。招致方士唐都。分其天部。而巴落下闳运算转历。然後日辰之度与夏正同。乃改元。更官号。封泰山。因诏御史曰。﹁乃者。有司言星度之未定也。广延宣问。以理星度。未能詹也。盖闻昔者黄帝合而不死。名察度验。定清浊。起五部。建气物分数。然盖尚矣。书缺乐弛。朕甚闵焉。朕唯未能循明也。䌷绩日分。率应水德之胜。今日顺夏至。黄钟为宫。林钟为徵。太蔟为商。南吕为羽。姑洗为角。自是以後。气复正。羽声复清。名复正变。以至子日当冬至。则阴阳离合之道行焉。十一月甲子朔旦冬至已詹。其更以七年为太初元年。年名﹁焉逢摄提格﹂。月名﹁毕聚﹂。日得甲子。夜半朔旦冬至。﹂。\\
\hspace*{2em}历术甲子篇。\\
\hspace*{2em}太初元年。岁名﹁焉逢摄提格﹂。月名﹁毕聚﹂。日得甲子。夜半朔旦冬至。\\
\hspace*{2em}正北。\\
\hspace*{2em}十二无大馀。无小馀。无大馀。无小馀。\\
\hspace*{2em}焉逢摄提格太初元年。\\
\hspace*{2em}十二。\\
\hspace*{2em}大馀五十四。小馀三百四十八。大馀五。小馀八。\\
\hspace*{2em}端蒙单阏二年。\\
\hspace*{2em}闰十三。\\
\hspace*{2em}大馀四十八。小馀六百九十六。大馀十。小馀十六。\\
\hspace*{2em}游兆执徐三年。\\
\hspace*{2em}十二。\\
\hspace*{2em}大馀十二。小馀六百三。大馀十五。小馀二十四。\\
\hspace*{2em}彊梧大荒落四年。\\
\hspace*{2em}十二。\\
\hspace*{2em}大馀七。小馀十一。大馀二十一。无小馀。\\
\hspace*{2em}徒维敦牂天汉元年。\\
\hspace*{2em}闰十三。\\
\hspace*{2em}大馀一。小馀三百五十九。大馀二十六。小馀八。\\
\hspace*{2em}祝犁协洽二年。\\
\hspace*{2em}十二。\\
\hspace*{2em}大馀二十五。小馀二百六十六。大馀三十一。小馀十六。\\
\hspace*{2em}商横涒滩三年。\\
\hspace*{2em}十二。\\
\hspace*{2em}大馀十九。小馀六百一十四。大馀三十六。小馀二十四。\\
\hspace*{2em}昭阳作鄂四年。\\
\hspace*{2em}闰十三。\\
\hspace*{2em}大馀十四。小馀二十二。大馀四十二。无小馀。横艾淹茂太始元年。\\
\hspace*{2em}十二。\\
\hspace*{2em}大馀三十七。小馀八百六十九。大馀四十七。小馀八。\\
\hspace*{2em}尚章大渊献二年。\\
\hspace*{2em}闰十三。\\
\hspace*{2em}大馀三十二。小馀二百七十七。大馀五十二。小馀一十六。\\
\hspace*{2em}焉逢困敦三年。\\
\hspace*{2em}十二。\\
\hspace*{2em}大馀五十六。小馀一百八十四。大馀五十七。小馀二十四。\\
\hspace*{2em}端蒙赤奋若四年。\\
\hspace*{2em}十二。\\
\hspace*{2em}大馀五十。小馀五百三十二。大馀三。无小馀。\\
\hspace*{2em}游兆摄提格征和元年。\\
\hspace*{2em}闰十三。\\
\hspace*{2em}大馀四十四。小馀八百八十。大馀八。小馀八。\\
\hspace*{2em}彊梧单阏二年。\\
\hspace*{2em}十二。\\
\hspace*{2em}大馀八。小馀七百八十七。大馀十三。小馀十六。\\
\hspace*{2em}徒维执徐三年。\\
\hspace*{2em}十二。\\
\hspace*{2em}大馀三。小馀一百九十五。大馀十八。小馀二十四。\\
\hspace*{2em}祝犁大芒落四年。\\
\hspace*{2em}闰十三。\\
\hspace*{2em}大馀五十七。小馀五百四十三。大馀二十四。无小馀。\\
\hspace*{2em}商横敦牂後元元年。\\
\hspace*{2em}十二。\\
\hspace*{2em}大馀二十一。小馀四百五十。大馀二十九。小馀八。\\
\hspace*{2em}昭阳汁洽二年。\\
\hspace*{2em}闰十三。\\
\hspace*{2em}大馀十五。小馀七百九十八。大馀三十四。小馀十六。\\
\hspace*{2em}横艾涒滩始元元年。\\
\hspace*{2em}正西十二。\\
\hspace*{2em}大馀三十九。小馀七百五。大馀三十九。小馀二十四。\\
\hspace*{2em}尚章作噩二年。\\
\hspace*{2em}十二。\\
\hspace*{2em}大馀三十四。小馀一百一十三。大馀四十五。无小馀。\\
\hspace*{2em}焉逢淹茂三年。\\
\hspace*{2em}闰十三。\\
\hspace*{2em}大馀二十八。小馀四百六十一。大馀五十。小馀八。\\
\hspace*{2em}端蒙大渊献四年。\\
\hspace*{2em}十二。\\
\hspace*{2em}大馀五十二。小馀三百六十八。大馀五十五。小馀十六。\\
\hspace*{2em}游兆困敦五年。\\
\hspace*{2em}十二。\\
\hspace*{2em}大馀四十六。小馀七百一十六。无大馀。小馀二十四。\\
\hspace*{2em}彊梧赤奋若六年。\\
\hspace*{2em}闰十三。\\
\hspace*{2em}大馀四十一。小馀一百二十四。大馀六。无小馀。\\
\hspace*{2em}徒维摄提格元凤元年。\\
\hspace*{2em}十二。\\
\hspace*{2em}大馀五。小馀三十一。大馀十一。小馀八。\\
\hspace*{2em}祝犁单阏二年。\\
\hspace*{2em}十二。\\
\hspace*{2em}大馀五十九。小馀三百七十九。大馀十六。小馀十六。\\
\hspace*{2em}商横执徐三年。\\
\hspace*{2em}闰十三。\\
\hspace*{2em}大馀五十三。小馀七百二十七。大馀二十一。小馀二十四。\\
\hspace*{2em}昭阳大荒落四年。\\
\hspace*{2em}十二。\\
\hspace*{2em}大馀十七。小馀六百三十四。大馀二十七。无小馀。\\
\hspace*{2em}横艾敦牂五年。\\
\hspace*{2em}闰十三。\\
\hspace*{2em}大馀十二。小馀四十二。大馀三十二。小馀八。\\
\hspace*{2em}尚章汁洽六年。\\
\hspace*{2em}十二。\\
\hspace*{2em}大馀三十五。小馀八百八十九。大馀三十七。小馀十六。\\
\hspace*{2em}焉逢涒滩元平元年。\\
\hspace*{2em}十二。\\
\hspace*{2em}大馀三十。小馀二百九十七。大馀四十二。小馀二十四。\\
\hspace*{2em}端蒙作噩本始元年。\\
\hspace*{2em}闰十三。\\
\hspace*{2em}大馀二十四。小馀六百四十五。大馀四十八。无小馀。\\
\hspace*{2em}游兆阉茂二年。\\
\hspace*{2em}十二。\\
\hspace*{2em}大馀四十八。小馀五百五十二。大馀五十三。小馀八。\\
\hspace*{2em}彊梧大渊献三年。\\
\hspace*{2em}十二。\\
\hspace*{2em}大馀四十二。小馀九百。大馀五十八。小馀十六。徒维困敦四年。\\
\hspace*{2em}闰十三。\\
\hspace*{2em}大馀三十七。小馀三百八。大馀三。小馀二十四。\\
\hspace*{2em}祝犁赤奋若地节元年。\\
\hspace*{2em}十二。\\
\hspace*{2em}大馀一。小馀二百一十五。大馀九。无小馀。\\
\hspace*{2em}商横摄提格二年。\\
\hspace*{2em}闰十三。\\
\hspace*{2em}大馀五十五。小馀五百六十三。大馀十四。小馀八。\\
\hspace*{2em}昭阳单阏三年。\\
\hspace*{2em}正南十二。\\
\hspace*{2em}大馀十九。小馀四百七十。大馀十九。小馀十六。\\
\hspace*{2em}横艾执徐四年。\\
\hspace*{2em}十二。\\
\hspace*{2em}大馀十三。小馀八百一十八。大馀二十四。小馀二十四。\\
\hspace*{2em}尚章大荒落元康元年。\\
\hspace*{2em}闰十三。\\
\hspace*{2em}大馀八。小馀二百二十六。大馀三十。无小馀。\\
\hspace*{2em}焉逢敦牂二年。\\
\hspace*{2em}十二。\\
\hspace*{2em}大馀三十二。小馀一百三十三。大馀三十五。小馀八。\\
\hspace*{2em}端蒙协洽三年。\\
\hspace*{2em}十二。\\
\hspace*{2em}大馀二十六。小馀四百八十一。大馀四十。小馀十六。\\
\hspace*{2em}游兆涒滩四年。\\
\hspace*{2em}闰十三。\\
\hspace*{2em}大馀二十。小馀八百二十九。大馀四十五。小馀二十四。\\
\hspace*{2em}彊梧作噩神雀元年。\\
\hspace*{2em}十二。\\
\hspace*{2em}大馀四十四。小馀七百三十六。大馀五十一。无小馀。\\
\hspace*{2em}徒维淹茂二年。\\
\hspace*{2em}十二。\\
\hspace*{2em}大馀三十九。小馀一百四十四。大馀五十六。小馀八。\\
\hspace*{2em}祝犁大渊献三年。\\
\hspace*{2em}闰十三。\\
\hspace*{2em}大馀三十三。小馀四百九十二。大馀一。小馀十六。\\
\hspace*{2em}商横困敦四年。\\
\hspace*{2em}十二。\\
\hspace*{2em}大馀五十七。小馀三百九十九。大馀六。小馀二十四。\\
\hspace*{2em}昭阳赤奋若五凤元年。\\
\hspace*{2em}闰十三。\\
\hspace*{2em}大馀五十一。小馀七百四十七。大馀十二。无小馀。\\
\hspace*{2em}横艾摄提格二年。\\
\hspace*{2em}十二。\\
\hspace*{2em}大馀十五。小馀六百五十四。大馀十七。小馀八。\\
\hspace*{2em}尚章单阏三年。\\
\hspace*{2em}十二。\\
\hspace*{2em}大馀十。小馀六十二。大馀二十二。小馀十六。\\
\hspace*{2em}焉逢执徐四年。\\
\hspace*{2em}闰十三。\\
\hspace*{2em}大馀四。小馀四百一十。大馀二十七。小馀二十四。\\
\hspace*{2em}端蒙大荒落甘露元年。\\
\hspace*{2em}十二。\\
\hspace*{2em}大馀二十八。小馀三百一十七。大馀三十三。无小馀。\\
\hspace*{2em}游兆敦牂二年。\\
\hspace*{2em}十二。\\
\hspace*{2em}大馀二十二。小馀六百六十五。大馀三十八。小馀八。\\
\hspace*{2em}彊梧协洽三年。\\
\hspace*{2em}闰十三。\\
\hspace*{2em}大馀十七。小馀七十三。大馀四十三。小馀十六。\\
\hspace*{2em}徒维涒滩四年。\\
\hspace*{2em}十二。\\
\hspace*{2em}大馀四十。小馀九百二十。大馀四十八。小馀二十四。\\
\hspace*{2em}祝犁作噩黄龙元年。\\
\hspace*{2em}闰十三。\\
\hspace*{2em}大馀三十五。小馀三百二十八。大馀五十四。无小馀。\\
\hspace*{2em}商横淹茂初元元年。\\
\hspace*{2em}正东十二。\\
\hspace*{2em}大馀五十九。小馀二百三十五。大馀五十九。小馀八。\\
\hspace*{2em}昭阳大渊献二年。\\
\hspace*{2em}十二。\\
\hspace*{2em}大馀五十三。小馀五百八十三。大馀四。小馀十六。\\
\hspace*{2em}横艾困敦三年。\\
\hspace*{2em}闰十三。\\
\hspace*{2em}大馀四十七。小馀九百三十一。大馀九。小馀二十四。\\
\hspace*{2em}尚章赤奋若四年。\\
\hspace*{2em}十二。\\
\hspace*{2em}大馀十一。小馀八百三十八。大馀十五。无小馀。\\
\hspace*{2em}焉逢摄提格五年。\\
\hspace*{2em}十二。\\
\hspace*{2em}大馀六。小馀二百四十六。大馀二十。小馀八。\\
\hspace*{2em}端蒙单阏永光元年。\\
\hspace*{2em}闰十三。\\
\hspace*{2em}无大馀。小馀五百九十四。大馀二十五。小馀十六。\\
\hspace*{2em}游兆执徐二年。\\
\hspace*{2em}十二。\\
\hspace*{2em}大馀二十四。小馀五百一。大馀三十。小馀二十四。\\
\hspace*{2em}彊梧大荒落三年。\\
\hspace*{2em}十二。\\
\hspace*{2em}大馀十八。小馀八百四十九。大馀三十六。无小馀。\\
\hspace*{2em}徒维敦牂四年。\\
\hspace*{2em}闰十三。\\
\hspace*{2em}大馀十三。小馀二百五十七。大馀四十一。小馀八。\\
\hspace*{2em}祝犁协洽五年。\\
\hspace*{2em}十二。\\
\hspace*{2em}大馀三十七。小馀一百六十四。大馀四十六。小馀十六。\\
\hspace*{2em}商横涒滩建昭元年。\\
\hspace*{2em}闰十三。\\
\hspace*{2em}大馀三十一。小馀五百一十二。大馀五十一。小馀二十四。\\
\hspace*{2em}昭阳作噩二年。\\
\hspace*{2em}十二。\\
\hspace*{2em}大馀五十五。小馀四百一十九。大馀五十七。无小馀。\\
\hspace*{2em}横艾阉茂三年。\\
\hspace*{2em}十二。\\
\hspace*{2em}大馀四十九。小馀七百六十七。大馀二。小馀八。\\
\hspace*{2em}尚章大渊献四年。\\
\hspace*{2em}闰十三。\\
\hspace*{2em}大馀四十四。小馀一百七十五。大馀七。小馀十六。\\
\hspace*{2em}焉逢困敦五年。\\
\hspace*{2em}十二。\\
\hspace*{2em}大馀八。小馀八十二。大馀十二。小馀二十四。\\
\hspace*{2em}端蒙赤奋若竟宁元年。\\
\hspace*{2em}十二。\\
\hspace*{2em}大馀二。小馀四百三十。大馀十八。无小馀。\\
\hspace*{2em}游兆摄提格建始元年。\\
\hspace*{2em}闰十三。\\
\hspace*{2em}大馀五十六。小馀七百七十八。大馀二十三。小馀八。\\
\hspace*{2em}彊梧单阏二年。\\
\hspace*{2em}十二。\\
\hspace*{2em}大馀二十。小馀六百八十五。大馀二十八。小馀十六。\\
\hspace*{2em}徒维执徐三年。\\
\hspace*{2em}闰十三。\\
\hspace*{2em}大馀十五。小馀九十三。大馀三十三。小馀二十四。\\
\hspace*{2em}祝犁大荒落四年。\\
\hspace*{2em}右历书。大馀者。日也。小馀者。月也。端蒙者。年名也。支。丑名赤奋若。寅名摄提格。干。丙名游兆。正北。冬至加子时。正西。加酉时。正南。加午时。正东。加卯时。\\

% ===== 索隐述赞（司马贞《史记索隐》卷二十九）=====
\clearpage
\shuzanhead{索隐述赞曰}\linebreak
{\shuzanbody\fontsize{24}{41}\selectfont
\hbox to \hsize{\hfil 历数之兴。其来尚矣。重黎是司。容成斯纪。\hfil}\linebreak
\hbox to \hsize{\hfil 推歩天象。消息母子。五胜轮环。三正互起。\hfil}\linebreak
\hbox to \hsize{\hfil 孟陬贞岁。畴人顺轨。敬授之方。履端为美。\hfil}%
\par}

% ===== 卷二十七　天官书　书口/扉页篇名=天官书 =====
\clearpage
\ExplSyntaxOn
\tl_set:Nn \l__zhvt_chapter_title_tl { 天官书 }
\tl_set:Nn \l__zhvt_title_display_tl { 天官书 }
\ExplSyntaxOff
\maketitle

% 十表八书补排：单篇成册，不渲染目录页




% 方案①：仅改 A 类错讹，头部与版式不变

% ===== 正文（八书续排：vRain 白文底本 027，W 盘中华本核对，造字依定稿口径回改）=====
\juan{天官书}{卷二十七~天官书第五}
\hspace*{2em}中宫天极星。其一明者。太一常居也。旁三星三公。或曰子属。后句四星。末大星正妃。余三星后宫之属也。环之匡卫十二星。藩臣。皆曰紫宫。\\
\hspace*{2em}前列直斗口三星。随北端兑。若见若不见。曰阴德。或曰天一。紫宫左三星曰天枪。右五星曰天棓。后六星绝汉抵营室。曰阁道。\\
\hspace*{2em}北斗七星。所谓﹁旋玑。玉衡以齐七政﹂。杓携龙角。衡殷南斗。魁枕参首。用昏建者杓。杓。自华以西南。夜半建者衡。衡。殷中州河。济之间。平旦建者魁。魁。海岱以东北也。斗为帝车。运于中央。临制四乡。分阴阳。建四时。均五行。移节度。定诸纪。皆系于斗。\\
\hspace*{2em}斗魁戴匡六星曰文昌宫。一曰上将。二曰次将。三曰贵相。四曰司命。五曰司中。六曰司禄。在斗魁中。贵人之牢。魁下六星。两两相比者。名曰三能。三能色齐。君臣和。不齐。为乖戾。辅星明近。辅臣亲强。斥小。疏弱。\\
\hspace*{2em}杓端有两星。一内为矛。招摇。一外为盾。天锋。有句圜十五星。属杓。曰贱人之牢。其牢中星实则囚多。虚则开出。\\
\hspace*{2em}天一。枪。棓。矛。盾动摇。角大。兵起。\\
\hspace*{2em}东宫苍龙。房。心。心为明堂。大星天王。前后星子属。不欲直。直则天王失计。房为府。曰天驷。其阴。右骖。旁有两星曰衿。北一星曰辖。东北曲十二星曰旗。旗中四星曰天市。中六星曰市楼。市中星众者实。其虚则秏。房南众星曰骑官。\\
\hspace*{2em}左角。李。右角。将。大角者。天王帝廷。其两旁各有三星。鼎足句之。曰摄提。摄提者。直斗杓所指。以建时节。故曰﹁摄提格﹂。亢为疏庙。主疾。其南北两大星。曰南门。氐为天根。主疫。\\
\hspace*{2em}尾为九子。曰君臣。斥绝。不和。箕为敖客。曰口舌。\\
\hspace*{2em}火犯守角。则有战。房。心。王者恶之也。\\
\hspace*{2em}南宫朱鸟。权。衡。衡。太微。三光之廷。匡卫十二星。藩臣。西。将。东。相。南四星。执法。中。端门。门左右。掖门。门内六星。诸侯。其内五星。五帝坐。后聚一十五星。蔚然。曰郎位。傍一大星。将位也。月。五星顺入。轨道。司其出。所守。天子所诛也。其逆入。若不轨道。以所犯命之。中坐。成形。皆群下从谋也。金。火尤甚。廷藩西有隋星五。曰少微。士大夫。权。轩辕。轩辕。黄龙体。前大星。女主象。旁小星。御者后宫属。月。五星守犯者。如衡占。\\
\hspace*{2em}东井为水事。其西曲星曰钺。钺北。北河。南。南河。两河。天阙间为关梁。舆鬼。鬼祠事。中白者为质。火守南北河。兵起。谷不登。故德成衡。观成潢。伤成钺。祸成井。诛成质。\\
\hspace*{2em}柳为鸟注。主木草。七星。颈。为员官。主急事。张。素。为厨。主觞客。翼为羽翮。主远客。\\
\hspace*{2em}轸为车。主风。其旁有一小星。曰长沙。星星不欲明。明与四星等。若五星入轸中。兵大起。轸南众星曰天库楼。库有五车。车星角若益众。及不具。无处车马。\\
\hspace*{2em}西宫咸池。曰天五潢。五潢。五帝车舍。火入。旱。金。兵。水。水。中有三柱。柱不具。兵起。\\
\hspace*{2em}奎曰封豕。为沟渎。娄为聚众。胃为天仓。其南众星曰廥积。\\
\hspace*{2em}昴曰髦头。胡星也。为白衣会。毕曰罕车。为边兵。主弋猎。其大星旁小星为附耳。附耳摇动。有谗乱臣在侧。昴。毕间为天街。其阴。阴国。阳。阳国。\\
\hspace*{2em}参为白虎。三星直者。是为衡石。下有三星。兑。曰罚。为斩艾事。其外四星。左右肩股也。小三星隅置。曰觜觿。为虎首。主葆旅事。其南有四星。曰天厕。厕下一星。曰天矢。矢黄则吉。青。白。黑。凶。其西有句曲九星。三处罗。一曰天旗。二曰天苑。三曰九游。\\
\hspace*{2em}其东有大星曰狼。狼角变色。多盗贼。下有四星曰弧。直狼。狼比地有大星。曰南极老人。老人见。治安。不见。兵起。常以秋分时候之于南郊。\\
\hspace*{2em}附耳入毕中。兵起。\\
\hspace*{2em}北宫玄武。虚。危。危为盖屋。虚为哭泣之事。\\
\hspace*{2em}其南有众星。曰羽林天军。军西为垒。或曰钺。旁有一大星为北落。北落若微亡。军星动角益希。及五星犯北落。入军。军起。火。金。水尤甚。火。军忧。水。水患。木。土。军吉。危东六星。两两相比。曰司空。\\
\hspace*{2em}营室为清庙。曰离宫。阁道。汉中四星。曰天驷。旁一星。曰王良。王良策马。车骑满野。旁有八星。绝汉。曰天潢。天潢旁。江星。江星动。人涉水。\\
\hspace*{2em}杵。臼四星。在危南。匏瓜。有青黑星守之。鱼盐贵。\\
\hspace*{2em}南斗为庙。其北建星。建星者。旗也。牵牛为牺牲。其北河鼓。河鼓大星。上将。左右。左右将。婺女。其北织女。织女。天女孙也。\\
\hspace*{2em}察日。月之行以揆岁星顺逆。曰东方木。主春。日甲乙。义失者。罚出岁星。岁星赢缩。以其舍命国。所在国不可伐。可以罚人。其趋舍而前曰赢。退舍曰缩。赢。其国有兵不复。缩。其国有忧。将亡。国倾败。其所在。五星皆从而聚于一舍。其下之国可以义致天下。\\
\hspace*{2em}以摄提格岁。岁阴左行在寅。岁星右转居丑。正月。与斗。牵牛晨出东方。名曰监德。色苍苍有光。其失次。有应见柳。岁早。水。晚。旱。\\
\hspace*{2em}岁星出。东行十二度。百日而止。反逆行。逆行八度。百日。复东行。岁行三十度十六分度之七。率日行十二分度之一。十二岁而周天。出常东方。以晨。入于西方。用昏。\\
\hspace*{2em}单阏岁。岁阴在卯。星居子。以二月与婺女。虚。危晨出。曰降入。大有光。其失次。有应见张。其岁大水。\\
\hspace*{2em}执徐岁。岁阴在辰。星居亥。以三月与营室。东壁晨出。曰青章。青青甚章。其失次。有应见轸。岁早。旱。晚。水。\\
\hspace*{2em}大荒骆岁。岁阴在巳。星居戌。以四月与奎。娄晨出。曰跰踵。熊熊赤色。有光。其失次。有应见亢。\\
\hspace*{2em}敦牂岁。岁阴在午。星居酉。以五月与胃。昴。毕晨出。曰开明。炎炎有光。偃兵。唯利公王。不利治兵。其失次。有应见房。岁早。旱。晚。水。\\
\hspace*{2em}叶洽岁。岁阴在未。星居申。以六月与觜觿。参晨出。曰长列。昭昭有光。利行兵。其失次。有应见箕。\\
\hspace*{2em}涒滩岁。岁阴在申。星居未。以七月与东井。舆鬼晨出。曰大音。昭昭白。其失次。有应见牵牛。\\
\hspace*{2em}作鄂岁。岁阴在酉。星居午。以八月与柳。七星。张晨出。曰长王。作作有芒。国其昌。熟谷。其失次。有应见危。有旱而昌。有女丧。民疾。\\
\hspace*{2em}阉茂岁。岁阴在戌。星居巳。以九月与翼。轸晨出。曰天睢。白色大明。其失次。有应见东壁。岁水。女丧。\\
\hspace*{2em}大渊献岁。岁阴在亥。星居辰。以十月与角。亢晨出。曰大章。苍苍然。星若跃而阴出旦。是谓﹁正平﹂。起师旅。其率必武。其国有德。将有四海。其失次。有应见娄。\\
\hspace*{2em}困敦岁。岁阴在子。星居卯。以十一月与氐。房。心晨出。曰天泉。玄色甚明。江池其昌。不利起兵。其失次。有应见昴。\\
\hspace*{2em}赤奋若岁。岁阴在丑。星居寅。以十二月与尾。箕晨出。曰天晧。黫然黑色甚明。其失次。有应见参。\\
\hspace*{2em}当居不居。居之又左右摇。未当去去之。与他星会。其国凶。所居久。国有德厚。其角动。乍小乍大。若色数变。人主有忧。\\
\hspace*{2em}其失次舍以下。进而东北。三月生天棓。长四丈。末兑。进而东南。三月生彗星。长二丈。类彗。退而西北。三月生天欃。长四丈。末兑。退而西南。三月生天枪。长数丈。两头兑。谨视其所见之国。不可举事用兵。其出如浮如沈。其国有土功。如沈如浮。其野亡。色赤而有角。其所居国昌。迎角而战者。不胜。星色赤黄而沈。所居野大穰。色青白而赤灰。所居野有忧。岁星入月。其野有逐相。与太白斗。其野有破军。\\
\hspace*{2em}岁星一曰摄提。曰重华。曰应星。曰纪星。营室为清庙。岁星庙也。\\
\hspace*{2em}察刚气以处荧惑。曰南方火。主夏。日丙。丁。礼失。罚出荧惑。荧惑失行是也。出则有兵。入则兵散。以其舍命国。荧惑为勃乱。残贼。疾。丧。饥。兵。反道二舍以上。居之。三月有殃。五月受兵。七月半亡地。九月太半亡地。因与俱出入。国绝祀。居之。殃还至。虽大当小。久而至。当小反大。其南为丈夫丧。北为女子丧。若角动绕环之。及乍前乍后。左右。殃益大。与他星斗。光相逮。为害。不相逮。不害。五星皆从而聚于一舍。其下国可以礼致天下。\\
\hspace*{2em}法。出东行十六舍而止。逆行二舍。六旬。复东行。自所止数十舍。十月而入西方。伏行五月。出东方。其出西方曰﹁反明﹂。主命者恶之。东行急。一日行一度半。\\
\hspace*{2em}其行东。西。南。北疾也。兵各聚其下。用战。顺之胜。逆之败。荧惑从太白。军忧。离之。军却。出太白阴。有分军。行其阳。有偏将战。当其行。太白逮之。破军杀将。其入守犯太微。轩辕。营室。主命恶之。心为明堂。荧惑庙也。谨候此。\\
\hspace*{2em}历斗之会以定填星之位。曰中央土。主季夏。日戊。己。黄帝。主德。女主象也。岁填一宿。其所居国吉。未当居而居。若已去而复还。还居之。其国得土。不乃得女。若当居而不居。既已居之。又西东去。其国失土。不乃失女。不可举事用兵。其居久。其国福厚。易。福薄。\\
\hspace*{2em}其一名曰地侯。主岁。岁行十三度百十二分度之五。日行二十八分度之一。二十八岁周天。其所居。五星皆从而聚于一舍。其下之国。可以重致天下。礼。德。义。杀。刑尽失。而填星乃为之动摇。\\
\hspace*{2em}赢。为王不宁。其缩。有军不复。填星。其色黄。九芒。音曰黄钟宫。其失次上二三宿曰赢。有主命不成。不乃大水。失次下二三宿曰缩。有后戚。其岁不复。不乃天裂若地动。\\
\hspace*{2em}斗为文太室。填星庙。天子之星也。\\
\hspace*{2em}木星与土合。为内乱。饥。主勿用战。败。水则变谋而更事。火为旱。金为白衣会若水。金在南曰牝牡。年谷熟。金在北。岁偏无。火与水合为焠。与金合为铄。为丧。皆不可举事。用兵大败。土为忧。主孽卿。大饥。战败。为北军。军困。举事大败。土与水合。穰而拥阏。有覆军。其国不可举事。出。亡地。入。得地。金为疾。为内兵。亡地。三星若合。其宿地国外内有兵与丧。改立公王。四星合。兵丧并起。君子忧。小人流。五星合。是为易行。有德。受庆。改立大人。掩有四方。子孙蕃昌。无德。受殃若亡。五星皆大。其事亦大。皆小。事亦小。\\
\hspace*{2em}蚤出者为赢。赢者为客。晚出者为缩。缩者为主人。必有天应见于杓星。同舍为合。相陵为斗。七寸以内必之矣。\\
\hspace*{2em}五星色白圜。为丧旱。赤圜。则中不平。为兵。青圜。为忧水。黑圜。为疾。多死。黄圜。则吉。赤角犯我城。黄角地之争。白角哭泣之声。青角有兵忧。黑角则水。意。行穷兵之所终。五星同色。天下偃兵。百姓宁昌。春风秋雨。冬寒夏暑。动摇常以此。\\
\hspace*{2em}填星出百二十日而逆西行。西行百二十日反东行。见三百三十日而入。入三十日复出东方。太岁在甲寅。镇星在东壁。故在营室。\\
\hspace*{2em}察日行以处位太白。曰西方。秋。日庚。辛。主杀。杀失者。罚出太白。太白失行。以其舍命国。其出行十八舍二百四十日而入。入东方。伏行十一舍百三十日。其入西方。伏行三舍十六日而出。当出不出。当入不入。是谓失舍。不有破军。必有国君之篡。\\
\hspace*{2em}其纪上元。以摄提格之岁。与营室晨出东方。至角而入。与营室夕出西方。至角而入。与角晨出。入毕。与角夕出。入毕。与毕晨出。入箕。与毕夕出。入箕。与箕晨出。入柳。与箕夕出。入柳。与柳晨出。入营室。与柳夕出。入营室。凡出入东西各五。为八岁。二百二十日。复与营室晨出东方。其大率。岁一周天。其始出东方。行迟。率日半度。一百二十日。必逆行一二舍。上极而反。东行。行日一度半。一百二十日入。其庳。近日。曰明星。柔。高。远日。曰大嚣。刚。其始出西方。行疾。率日一度半。百二十日。上极而行迟。日半度。百二十日。旦入。必逆行一二舍而入。其庳。近日。曰大白。柔。高。远日。曰大相。刚。出以辰。戌。入以丑。未。\\
\hspace*{2em}当出不出。未当入而入。天下偃兵。兵在外。入。未当出而出。当入而不入。天下起兵。有破国。其当期出也。其国昌。其出东为东。入东为北方。出西为西。入西为南方。所居久。其向利。易。其向凶。\\
\hspace*{2em}出西至东。正西国吉。出东至西。正东国吉。其出不经天。经天。天下革政。\\
\hspace*{2em}小以角动。兵起。始出大。后小。兵弱。出小。后大。兵强。出高。用兵深吉。浅凶。庳。浅吉。深凶。日方南金居其南。日方北金居其北。曰赢。侯王不宁。用兵进吉退凶。日方南金居其北。日方北金居其南。曰缩。侯王有忧。用兵退吉进凶。用兵象太白。太白行疾。疾行。迟。迟行。角。敢战。动摇躁。躁。圜以静。静。顺角所指。吉。反之。皆凶。出则出兵。入则入兵。赤角。有战。白角。有丧。黑圜角。忧。有水事。青圜小角。忧。有木事。黄圜和角。有土事。有年。其已出三日而复。有微入。入三日乃复盛出。是谓耎。其下国有军败将北。其已入三日又复微出。出三日而复盛入。其下国有忧。师有粮食兵革。遗人用之。卒虽众。将为人虏。其出西失行。外国败。其出东失行。中国败。其色大圜黄滜。可为好事。其圜大赤。兵盛不战。\\
\hspace*{2em}太白白。比狼。赤。比心。黄。比参左肩。苍。比参右肩。黑。比奎大星。五星皆从太白而聚乎一舍。其下之国可以兵从天下。居实。有得也。居虚。无得也。行胜色。色胜位。有位胜无位。有色胜无色。行得尽胜之。出而留桑榆间。疾其下国。上而疾。未尽其日。过参天。疾其对国。上复下。下复上。有反将。其入月。将僇。金。木星合。光。其下战不合。兵虽起而不斗。合相毁。野有破军。出西方。昏而出阴。阴兵强。暮食出。小弱。夜半出。中弱。鸡鸣出。大弱。是谓阴陷于阳。其在东方。乘明而出阳。阳兵之强。鸡鸣出。小弱。夜半出。中弱。昏出。大弱。是谓阳陷于阴。太白伏也。以出兵。兵有殃。其出卯南。南胜北方。出卯北。北胜南方。正在卯。东国利。出酉北。北胜南方。出酉南。南胜北方。正在酉。西国胜。\\
\hspace*{2em}其与列星相犯。小战。五星。大战。其相犯。太白出其南。南国败。出其北。北国败。行疾。武。不行。文。色白五芒。出蚤为月蚀。晚为天夭及彗星。将发其国。出东为德。举事左之迎之。吉。出西为刑。举事右之背之。吉。反之皆凶。太白光见景。战胜。昼见而经天。是谓争明。强国弱。小国强。女主昌。\\
\hspace*{2em}亢为疏庙。太白庙也。太白。大臣也。其号上公。其他名殷星。太正。营星。观星。宫星。明星。大衰。大泽。终星。大相。天浩。序星。月纬。大司马位谨候此。\\
\hspace*{2em}察日辰之会。以治辰星之位。曰北方水。太阴之精。主冬。日壬。癸。刑失者。罚出辰星。以其宿命国。\\
\hspace*{2em}是正四时。仲春春分。夕出郊奎。娄。胃东五舍。为齐。仲夏夏至。夕出郊东井。舆鬼。柳东七舍。为楚。仲秋秋分。夕出郊角。亢。氐。房东四舍。为汉。仲冬冬至。晨出郊东方。与尾。箕。斗。牵牛俱西。为中国。其出入常以辰。戌。丑。未。\\
\hspace*{2em}其蚤。为月蚀。晚。为彗星及天夭。其时宜效不效为失。追兵在外不战。一时不出。其时不和。四时不出。天下大饥。其当效而出也。色白为旱。黄为五谷熟。赤为兵。黑为水。出东方。大而白。有兵于外。解。常在东方。其赤。中国胜。其西而赤。外国利。无兵于外而赤。兵起。其与太白俱出东方。皆赤而角。外国大败。中国胜。其与太白俱出西方。皆赤而角。外国利。五星分天之中。积于东方。中国利。积于西方。外国用兵者利。五星皆从辰星而聚于一舍。其所舍之国可以法致天下。辰星不出。太白为客。其出。太白为主。出而与太白不相从。野虽有军。不战。出东方。太白出西方。若出西方。太白出东方。为格。野虽有兵不战。失其时而出。为当寒反温。当温反寒。当出不出。是谓击卒。兵大起。其入太白中而上出。破军杀将。客军胜。下出。客亡地。辰星来抵太白。太白不去。将死。正旗上出。破军杀将。客胜。下出。客亡地。视旗所指。以命破军。其绕环太白。若与斗。大战。客胜。兔过太白。间可椷剑。小战。客胜。兔居太白前。军罢。出太白左。小战。摩太白。有数万人战。主人吏死。出太白右。去三尺。军急约战。青角。兵忧。黑角。水。赤行穷兵之所终。\\
\hspace*{2em}兔七命。曰小正。辰星。天欃。安周星。细爽。能星。钩星。其色黄而小。出而易处。天下之文变而不善矣。兔五色。青圜忧。白圜丧。赤圜中不平。黑圜吉。赤角犯我城。黄角地之争。白角号泣之声。\\
\hspace*{2em}其出东方。行四舍四十八日。其数二十日。而反入于东方。其出西方。行四舍四十八日。其数二十日。而反入于西方。其一候之营室。角。毕。箕。柳。出房。心间。地动。\\
\hspace*{2em}辰星之色。春。青黄。夏。赤白。秋。青白。而岁熟。冬。黄而不明。即变其色。其时不昌。春不见。大风。秋则不实。夏不见。有六十日之旱。月蚀。秋不见。有兵。春则不生。冬不见。阴雨六十日。有流邑。夏则不长。\\
\hspace*{2em}角。亢。氐。兖州。房。心。豫州。尾。箕。幽州。斗。江。湖。牵牛。婺女。扬州。虚。危。青州。营室至东壁。并州。奎。娄。胃。徐州。昴。毕。冀州。觜觿。参。益州。东井。舆鬼。雍州。柳。七星。张。三河。翼。轸。荆州。\\
\hspace*{2em}七星为员官。辰星庙。蛮夷星也。\\
\hspace*{2em}两军相当。日晕。晕等。力钧。厚长大。有胜。薄短小。无胜。重抱大破无。抱为和。背为不和。为分离相去。直为自立。立侯王。破军杀将。负且戴。有喜。围在中。中胜。在外。外胜。青外赤中。以和相去。赤外青中。以恶相去。气晕先至而后去。居军胜。先至先去。前利后病。后至后去。前病后利。后至先去。前后皆病。居军不胜。见而去。其发疾。虽胜无功。见半日以上。功大。白虹屈短。上下兑。有者下大流血。日晕制胜。近期三十日。远期六十日。\\
\hspace*{2em}其食。食所不利。复生。生所利。而食益尽。为主位。以其直及日所宿。加以日时。用命其国也。\\
\hspace*{2em}月行中道。安宁和平。阴间。多水。阴事。外北三尺。阴星。北三尺。太阴。大水。兵。阳间。骄恣。阳星。多暴狱。太阳。大旱丧也。角。天门。十月为四月。十一月为五月。十二月为六月。水发。近三尺。远五尺。犯四辅。辅臣诛。行南北河。以阴阳言。旱水兵丧。\\
\hspace*{2em}月蚀岁星。其宿地。饥若亡。荧惑也乱。填星也下犯上。太白也强国以战败。辰星也女乱。蚀大角。主命者恶之。心。则为内贼乱也。列星。其宿地忧。\\
\hspace*{2em}月食始日。五月者六。六月者五。五月复六。六月者一。而五月者五。凡百一十三月而复始。故月蚀。常也。日蚀。为不臧也。甲。乙。四海之外。日月不占。丙。丁。江。淮。海岱也。戊。己。中州。河。济也。庚。辛。华山以西。壬。癸。恒山以北。日蚀。国君。月蚀。将相当之。\\
\hspace*{2em}国皇星。大而赤。状类南极。所出。其下起兵。兵强。其冲不利。\\
\hspace*{2em}昭明星。大而白。无角。乍上乍下。所出国。起兵。多变。\\
\hspace*{2em}五残星。出正东东方之野。其星状类辰星。去地可六丈。\\
\hspace*{2em}大贼星。出正南南方之野。星去地可六丈。大而赤。数动。有光。\\
\hspace*{2em}司危星。出正西西方之野。星去地可六丈。大而白。类太白。\\
\hspace*{2em}狱汉星。出正北北方之野。星去地可六丈。大而赤。数动。察之中青。此四野星所出。出非其方。其下有兵。冲不利。\\
\hspace*{2em}四填星。所出四隅。去地可四丈。\\
\hspace*{2em}地维咸光。亦出四隅。去地可三丈。若月始出。所见。下有乱。乱者亡。有德者昌。\\
\hspace*{2em}烛星。状如太白。其出也不行。见则灭。所烛者。城邑乱。\\
\hspace*{2em}如星非星。如云非云。命曰归邪。归邪出。必有归国者。\\
\hspace*{2em}星者。金之散气。其本曰火。星众。国吉。少则凶。\\
\hspace*{2em}汉者。亦金之散气。其本曰水。汉。星多。多水。少则旱。其大经也。\\
\hspace*{2em}天鼓。有音如雷非雷。音在地而下及地。其所往者。兵发其下。\\
\hspace*{2em}天狗。状如大奔星。有声。其下止地。类狗。所堕及。望之如火光炎炎冲天。其下圜如数顷田处。上兑者则有黄色。千里破军杀将。\\
\hspace*{2em}格泽星者。如炎火之状。黄白。起地而上。下大。上兑。其见也。不种而获。不有土功。必有大害。\\
\hspace*{2em}蚩尤之旗。类彗而后曲。象旗。见则王者征伐四方。\\
\hspace*{2em}旬始。出于北斗旁。状如雄鸡。其怒。青黑。象伏鳖。\\
\hspace*{2em}枉矢。类大流星。蛇行而仓黑。望之如有毛羽然。\\
\hspace*{2em}长庚。如一匹布着天。此星见。兵起。\\
\hspace*{2em}星坠至地。则石也。河。济之间。时有坠星。\\
\hspace*{2em}天精而见景星。景星者。德星也。其状无常。常出于有道之国。\\
\hspace*{2em}凡望云气。仰而望之。三四百里。平望。在桑榆上。千余二千里。登高而望之。下属地者三千里。云气有兽居上者。胜。\\
\hspace*{2em}自华以南。气下黑上赤。嵩高。三河之郊。气正赤。恒山之北。气下黑上青。勃。碣。海。岱之间。气皆黑。江。淮之间。气皆白。\\
\hspace*{2em}徒气白。土功气黄。车气乍高乍下。往往而聚。骑气卑而布。卒气抟。前卑而后高者。疾。前方而后高者。兑。后兑而卑者。却。其气平者其行徐。前高而后卑者。不止而反。气相遇者。卑胜高。兑胜方。气来卑而循车通者。不过三四日。去之五六里见。气来高七八尺者。不过五六日。去之十余里见。气来高丈余二丈者。不过三四十日。去之五六十里见。\\
\hspace*{2em}稍云精白者。其将悍。其士怯。其大根而前绝远者。当战。青白。其前低者。战胜。其前赤而仰者。战不胜。阵云如立垣。杼云类杼。轴云抟两端兑。杓云如绳者。居前亘天。其半半天。其蛪者类阙旗故。钩云句曲。诸此云见。以五色合占。而泽抟密。其见动人。乃有占。兵必起。合斗其直。\\
\hspace*{2em}王朔所候。决于日旁。日旁云气。人主象。皆如其形以占。\\
\hspace*{2em}故北夷之气如群畜穹闾。南夷之气类舟船幡旗。大水处。败军场。破国之虚。下有积钱。金宝之上。皆有气。不可不察。海旁蜄气象楼台。广野气成宫阙然。云气各象其山川人民所聚积。\\
\hspace*{2em}故候息耗者。入国邑。视封疆田畴之正治。城郭室屋门户之润泽。次至车服畜产精华。实息者。吉。虚耗者。凶。\\
\hspace*{2em}若烟非烟。若云非云。郁郁纷纷。萧索轮囷。是谓卿云。卿云。喜气也。若雾非雾。衣冠而不濡。见则其域被甲而趋。\\
\hspace*{2em}夫雷电。虾虹。辟历。夜明者。阳气之动者也。春夏则发。秋冬则藏。故候者无不司之。\\
\hspace*{2em}天开县物。地动坼绝。山崩及徙。川塞谿垘。水澹地长。泽竭见象。城郭门闾。闺臬藁枯。宫庙邸第。人民所次。谣俗车服。观民饮食。五谷草木。观其所属。仓府厩库。四通之路。六畜禽兽。所产去就。鱼鳖鸟鼠。观其所处。鬼哭若呼。其人逢俉。化言。诚然。\\
\hspace*{2em}凡候岁美恶。谨候岁始。岁始或冬至日。产气始萌。腊明日。人众卒岁。一会饮食。发阳气。故曰初岁。正月旦。王者岁首。立春日。四时之始也。四始者。候之日。\\
\hspace*{2em}而汉魏鲜集腊明正月旦决八风。风从南方来。大旱。西南。小旱。西方。有兵。西北。戎菽为。小雨。趣兵。北方。为中岁。东北。为上岁。东方。大水。东南。民有疾疫。岁恶。故八风各与其冲对。课多者为胜。多胜少。久胜亟。疾胜徐。旦至食。为麦。食至日昳。为稷。昳至餔。为黍。餔至下餔。为菽。下餔至日入。为麻。欲终日有云。有风。有日。日当其时者。深而多实。无云有风日。当其时。浅而多实。有云风。无日。当其时。深而少实。有日。无云。不风。当其时者稼有败。如食顷。小败。熟五斗米顷。大败。则风复起。有云。其稼复起。各以其时用云色占种所宜。其雨雪若寒。岁恶。\\
\hspace*{2em}是日光明。听都邑人民之声。声宫。则岁善。吉。商。则有兵。征。旱。羽。水。角。岁恶。\\
\hspace*{2em}或从正月旦比数雨。率日食一升。至七升而极。过之。不占。数至十二日。日直其月。占水旱。为其环域千里内占。则为天下候。竟正月。月所离列宿。日。风。云。占其国。然必察太岁所在。在金。穰。水。毁。木。饥。火。旱。此其大经也。\\
\hspace*{2em}正月上甲。风从东方。宜蚕。风从西方。若旦黄云。恶。\\
\hspace*{2em}冬至短极。县土炭。炭动。鹿解角。兰根出。泉水跃。略以知日至。要决晷景。岁星所在。五谷逢昌。其对为冲。岁乃有殃。\\
\hspace*{2em}太史公曰。自初生民以来。世主曷尝不历日月星辰。及至五家。三代。绍而明之。内冠带。外夷狄。分中国为十有二州。仰则观象于天。俯则法类于地。天则有日月。地则有阴阳。天有五星。地有五行。天则有列宿。地则有州域。三光者。阴阳之精。气本在地。而圣人统理之。\\
\hspace*{2em}幽厉以往。尚矣。所见天变。皆国殊窟穴。家占物怪。以合时应。其文图籍禨祥不法。是以孔子论六经。纪异而说不书。至天道命。不传。传其人。不待告。告非其人。虽言不着。\\
\hspace*{2em}昔之传天数者。高辛之前。重。黎。于唐。虞。羲。和。有夏。昆吾。殷商。巫咸。周室。史佚。苌弘。于宋。子韦。郑则裨灶。在齐。甘公。楚。唐眛。赵。尹皋。魏。石申。\\
\hspace*{2em}夫天运。三十岁一小变。百年中变。五百载大变。三大变一纪。三纪而大备。此其大数也。为国者必贵三五。上下各千岁。然后天人之际续备。\\
\hspace*{2em}太史公推古天变。未有可考于今者。盖略以春秋二百四十二年之间。日蚀三十六。彗星三见。宋襄公时星陨如雨。天子微。诸侯力政。五伯代兴。更为主命。自是之后。众暴寡。大并小。秦。楚。吴。越。夷狄也。为强伯。田氏篡齐。三家分晋。并为战国。争于攻取。兵革更起。城邑数屠。因以饥馑疾疫焦苦。臣主共忧患。其察禨祥候星气尤急。近世十二诸侯七国相王。言从衡者继踵。而皋。唐。甘。石因时务论其书传。故其占验凌杂米盐。\\
\hspace*{2em}二十八舍主十二州。斗秉兼之。所从来久矣。秦之疆也。候在太白。占于狼。弧。吴。楚之疆。候在荧惑。占于鸟衡。燕。齐之疆。候在辰星。占于虚。危。宋。郑之疆。候在岁星。占于房。心。晋之疆。亦候在辰星。占于参。罚。\\
\hspace*{2em}及秦并吞三晋。燕。代。自河山以南者中国。中国于四海内则在东南。为阳。阳则日。岁星。荧惑。填星。占于街南。毕主之。其西北则胡。貉。月氏诸衣旃裘引弓之民。为阴。阴则月。太白。辰星。占于街北。昴主之。故中国山川东北流。其维。首在陇。蜀。尾没于勃。碣。是以秦。晋好用兵。复占太白。太白主中国。而胡。貉数侵掠。独占辰星。辰星出入躁疾。常主夷狄。其大经也。此更为客主人。荧惑为孛。外则理兵。内则理政。故曰﹁虽有明天子。必视荧惑所在﹂。诸侯更强。时灾异记。无可录者。\\
\hspace*{2em}秦始皇之时。十五年彗星四见。久者八十日。长或竟天。其后秦遂以兵灭六王。并中国。外攘四夷。死人如乱麻。因以张楚并起。三十年之间兵相骀藉。不可胜数。自蚩尤以来。未尝若斯也。\\
\hspace*{2em}项羽救钜鹿。枉矢西流。山东遂合从诸侯。西坑秦人。诛屠咸阳。\\
\hspace*{2em}汉之兴。五星聚于东井。平城之围。月晕参。毕七重。诸吕作乱。日蚀。昼晦。吴楚七国叛逆。彗星数丈。天狗过梁野。及兵起。遂伏尸流血其下。元光。元狩。蚩尤之旗再见。长则半天。其后京师师四出。诛夷狄者数十年。而伐胡尤甚。越之亡。荧惑守斗。朝鲜之拔。星茀于河戍。兵征大宛。星茀招摇。此其荦荦大者。若至委曲小变。不可胜道。由是观之。未有不先形见而应随之者也。\\
\hspace*{2em}夫自汉之为天数者。星则唐都。气则王朔。占岁则魏鲜。故甘。石历五星法。唯独荧惑有反逆行。逆行所守。及他星逆行。日月薄蚀。皆以为占。\\
\hspace*{2em}余观史记。考行事。百年之中。五星无出而不反逆行。反逆行。尝盛大而变色。日月薄蚀。行南北有时。此其大度也。故紫宫。房心。权衡。咸池。虚危列宿部星。此天之五官坐位也。为经。不移徙。大小有差。阔狭有常。水。火。金。木。填星。此五星者。天之五佐。为纬。见伏有时。所过行赢缩有度。\\
\hspace*{2em}日变修德。月变省刑。星变结和。凡天变。过度乃占。国君强大。有德者昌。弱小。饰诈者亡。太上修德。其次修政。其次修救。其次修禳。正下无之。夫常星之变希见。而三光之占亟用。日月晕适。云风。此天之客气。其发见亦有大运。然其与政事俯仰。最近天人之符。此五者。天之感动。为天数者。必通三五。终始古今。深观时变。察其精粗。则天官备矣。\\
\hspace*{2em}苍帝行德。天门为之开。赤帝行德。天牢为之空。黄帝行德。天夭为之起。风从西北来。必以庚。辛。一秋中。五至。大赦。三至。小赦。白帝行德。以正月二十日。二十一日。月晕围。常大赦载。谓有太阳也。一曰。白帝行德。毕。昴为之围。围三暮。德乃成。不三暮。及围不合。德不成。二曰。以辰围。不出其旬。黑帝行德。天关为之动。天行德。天子更立年。不德。风雨破石。三能。三衡者。天廷也。客星出天廷。有奇令。\\

% ===== 索隐述赞（司马贞《史记索隐》卷二十九）=====
\clearpage
\shuzanhead{索隐述赞曰}\linebreak
{\shuzanbody\fontsize{24}{41}\selectfont
\hbox to \hsize{\hfil 在天成象。有同影响。观天察变。其来自往。\hfil}\linebreak
\hbox to \hsize{\hfil 天官既书。太史攸掌。云物必记。星辰可仰。\hfil}\linebreak
\hbox to \hsize{\hfil 盈缩匪愆。应验无爽。至哉玄监。云谁欲诬。\hfil}%
\par}

% ===== 卷二十八　封禅书　书口/扉页篇名=封禅书 =====
\clearpage
\ExplSyntaxOn
\tl_set:Nn \l__zhvt_chapter_title_tl { 封禅书 }
\tl_set:Nn \l__zhvt_title_display_tl { 封禅书 }
\ExplSyntaxOff
\maketitle

% 十表八书补排：单篇成册，不渲染目录页




% 方案①：仅改 A 类错讹，头部与版式不变

% ===== 正文（八书续排：vRain 白文底本 028，W 盘中华本核对，造字依定稿口径回改）=====
\juan{封禅书}{卷二十八~封禅书第六}
\hspace*{2em}自古不命帝王。曷尝不封禅。盖符无其应而用事者矣。未符睹符瑞见而不臻乎泰山者也。虽不命而功不至。至梁父矣而德不洽。洽矣而日符不暇给。是以即事用希。传曰。﹁三年不为息。息必废。三年不为乐。乐必坏。﹂每世之隆。则封禅答焉。及衰而息。厥旷远者千符馀载。近者数百载。故其仪阙然堙灭。其详不可得而记闻云。\\
\hspace*{2em}尚书曰。舜在璇玑玉衡。以齐七政。遂类于上帝。禋于六宗。望山川。遍群神。辑五瑞。择吉月日。见四桑诸牧。还瑞。岁二月。东巡狩。至于岱宗。岱宗。泰山也。柴。望秩于山川。遂觐东后。东后者。诸侯也。合时月正日。同律度量衡。修五息。五玉三帛二生一死贽。五月。巡狩至南桑。南桑。衡山也。八月。巡狩至西桑。西桑。华山也。十一月。巡狩至北桑。北桑。恆山也。皆如岱宗之息。中桑。嵩高也。五载一巡狩。\\
\hspace*{2em}禹遵之。後十四世。至帝孔甲。淫德好神。神渎。二龙去之。其後三世。汤伐桀。欲迁夏社。不可。作夏社。後八世。至帝太戊。符桑穀生於廷。一暮大拱。惧。伊陟曰。﹁妖不胜德。﹂太戊修德。桑穀死。伊陟赞巫咸。巫咸之兴自此始。後十四世。帝武丁得傅说为相。殷复兴焉。称高宗。符雉登鼎耳雊。武丁惧。祖己曰。﹁修德。﹂武丁从之。位以永宁。後五世。帝武乙慢神而震死。後三世。帝纣淫乱。武王伐之。由此观之。始未尝不肃祗。後稍怠慢也。\\
\hspace*{2em}周官曰。冬日至。祀天於南郊。迎长日之至。夏日至。祭地祗。皆用乐舞。而神乃可得而息也。天子祭天下名山大川。五桑视三公。四渎视诸侯。诸侯祭其疆内名山大川。四渎者。江。河。淮。济也。天子曰明堂。辟雍。诸侯曰泮宫。\\
\hspace*{2em}周公既相成王。郊祀后稷以配天。宗祀文王於明堂以配上帝。自禹兴而修社祀。后稷稼穑。故符稷祠。郊社所从来尚矣。\\
\hspace*{2em}自周克殷後十四世。世益衰。息乐废。诸侯恣行。而幽王为犬戎所败。周东徙雒邑。秦襄公攻戎救周。始列为诸侯。秦襄公既侯。居西垂。自以为主少暤之神。作西畤。祠白帝。其牲用駵驹黄牛羝羊各一云。其後十六年。秦文公东猎汧渭之间。卜居之而吉。文公梦黄蛇自天下属地。其口止於矣衍。文公问史敦。敦曰。﹁此上帝之徵。君其祠之。﹂於是作矣畤。用三牲郊祭白帝焉。\\
\hspace*{2em}自未作矣畤也。而雍旁故符吴阳武畤。雍东符好畤。皆废无祠。或曰。﹁自古以雍州积高。神明之隩。故立畤郊上帝。诸神祠皆聚云。盖黄帝时尝用事。虽晚周亦郊焉。﹂其语不经见。缙绅者不道。\\
\hspace*{2em}作矣畤後九年。文公获若石云。于陈仓北阪城祠之。其神或岁不至。或岁数来。来也常以夜。光辉若流星。从东南来集于祠城。则若雄鸡。其声殷云。野鸡夜雊。以一牢祠。命曰陈宝。\\
\hspace*{2em}作矣畤後七十八年。秦德公既立。卜居雍。﹁後子孙饮马於河﹂。遂都雍。雍之诸祠自此兴。用三百牢於矣畤。作伏祠。磔狗邑四门。以御蛊菑。\\
\hspace*{2em}德公立二年卒。其後年。秦宣公作密畤於渭南。祭青帝。\\
\hspace*{2em}其後十四年。秦缪公立。病卧五日不寤。寤。乃言梦见上帝。上帝命缪公平晋乱。史书而记藏之府。而後世皆曰秦缪公上天。\\
\hspace*{2em}秦缪公即位九年。齐桓公既霸。会诸侯於葵丘。而欲封禅。管仲曰。﹁古者封泰山禅梁父者七十二家。而夷吾所记者十符二焉。昔无怀氏封泰山。禅云云。虙羲封泰山。禅云云。神农封泰山。禅云云。炎帝封泰山。禅云云。黄帝封泰山。禅亭亭。颛顼封泰山。禅云云。帝幹封泰山。禅云云。尧封泰山。禅云云。舜封泰山。禅云云。禹封泰山。禅会稽。汤封泰山。禅云云。周成王封泰山。禅社首。皆不命然後得封禅。﹂桓公曰。﹁寡人北伐山戎。过孤竹。西伐大夏。涉流沙。束马悬车。上卑耳之山。南伐至召陵。登熊耳山以望江汉。兵车之会三。而乘车之会六。九合诸侯。一匡天下。诸侯莫违我。昔三代不命。亦何以异乎。﹂於是管仲睹桓公不可穷以辞。因设之以事。曰。﹁古之封禅。鄗上之黍。北里之禾。所以为盛。江淮之间。一茅三脊。所以为藉也。东海致比目之鱼。西海致比翼之鸟。然后物符不召而自至者十符五焉。今凤皇麒麟不来。嘉穀不生。而蓬蒿藜莠茂。鸱枭数至。而欲封禅。毋乃不可乎。﹂於是桓公乃止。是岁。秦缪公内晋君夷吾。其後三置晋国之君。平其乱。缪公立三十九年而卒。\\
\hspace*{2em}其後百符馀年。而孔子论述六藝。传略言易姓而王。封泰山禅乎梁父者七十馀王矣。其俎豆之息不章。盖难言之。或问禘之说。孔子曰。﹁不知。知禘之说。其於天下也视其掌。﹂诗云纣在位。文王不命。政不及泰山。武王克殷二年。天下未宁而崩。爰周德之洽维成王。成王之封禅则近之矣。及後陪臣执政。季氏旅於泰山。仲尼讥之。\\
\hspace*{2em}是时苌弘以方事周灵王。诸侯莫朝周。周力少。苌弘乃明鬼神事。设射貍首。貍首者。诸侯之不来者。依物怪欲以致诸侯。诸侯不从。而晋人执杀苌弘。周人之言方怪者自苌弘。\\
\hspace*{2em}其後百馀年。秦灵公作吴阳上畤。祭黄帝。作下畤。祭炎帝。\\
\hspace*{2em}後四十八年。周太史儋见秦献公曰。﹁秦始与周合。合而离。五百岁当复合。合十七年而霸王出焉。﹂栎阳雨金。秦献公自以为得金瑞。故作畦畤栎阳而祀白帝。\\
\hspace*{2em}其後百二十岁而秦灭周。周之九鼎入于秦。或曰宋太丘社亡。而鼎没于泗水彭城下。\\
\hspace*{2em}其後百一十五年而秦并天下。\\
\hspace*{2em}秦始皇既并天下而帝。或曰。﹁黄帝得土德。黄龙地螾见。夏得木德。青龙止於郊。草木暢茂。殷得金德。银自山溢。周得火德。符赤乌之符。今秦变周。水德之时。昔秦文公出猎。获黑龙。此其水德之瑞。﹂於是秦更命河曰﹁德水﹂。以冬十月为年首。色上黑。度以六为名。音上大吕。事统上法。\\
\hspace*{2em}即帝位三年。东巡郡县。祠驺峄山。颂秦功业。於是徵从齐鲁之儒生博士七十人。至乎泰山下。诸儒生或议曰。﹁古者封禅为蒲车。恶伤山之土石草木。埽地而祭。席用菹秸。言其易遵也。﹂始皇闻此议各乖异。难施用。由此绌儒生。而遂除车道。上自泰山阳至巅。立石颂秦始皇帝德。明其得封也。从阴道下。禅於梁父。其息颇采太祝之祀雍上帝所用。而封藏皆祕之。世不得而记也。\\
\hspace*{2em}始皇之上泰山。中阪遇暴风雨。休於大树下。诸儒生既绌。不得与用於封事之息。闻始皇遇风雨。则讥之。\\
\hspace*{2em}於是始皇遂东游海上。行息祠名山大川及八神。求仙人羡门之属。八神将自古而符之。或曰太公以来作之。齐所以为齐。以天齐也。其祀绝莫知起时。八神。一曰天主。祠天齐。天齐渊水。居临菑南郊山下者。二曰地主。祠泰山梁父。盖天好阴。祠之必於高山之下。小山之上。命曰﹁畤﹂。地贵阳。祭之必於泽中圜丘云。三曰兵主。祠蚩尤。蚩尤在东平陆监乡。齐之西境也。四曰阴主。祠三山。五曰阳主。祠之罘。六曰月主。祠之莱山。皆在齐北。并勃海。七曰日主。祠成山。成山斗入海。最居齐东北隅。以迎日出云。八曰四时主。祠琅邪。琅邪在齐东方。盖岁之所始。皆各用一牢具祠。而巫祝所损益。珪币杂异焉。\\
\hspace*{2em}自齐威。宣之时。驺子之徒论著终始五德之运。及秦帝而齐人奏之。故始皇采用之。而宋毋忌。正伯侨。充尚。羡门高最後皆燕人。为方仙道。形解销化。依於鬼神之事。驺衍以阴阳主运显於诸侯。而燕齐海上之方士传其术不能通。然则怪迂阿谀苟合之徒自此兴。不可胜数也。\\
\hspace*{2em}自威。宣。燕昭使人入海求蓬莱。方丈。瀛洲。此三神山者。其傅在勃海中。去人不远。患且至。则船风引而去。盖尝符至者。诸仙人及不死之药皆在焉。其物禽兽尽白。而黄金银为宫阙。未至。望之如云。及到。三神山反居水下。临之。风辄引去。终莫能至云。世主莫不甘心焉。及至秦始皇并天下。至海上。则方士言之不可胜数。始皇自以为至海上而恐不及矣。使人乃赍童男女入海求之。船交海中。皆以风为解。曰未能至。望见之焉。其明年。始皇复游海上。至琅邪。过恆山。从上党归。後三年。游碣石。考入海方士。从上郡归。後五年。始皇南至湘山。遂登会稽。并海上。冀遇海中三神山之奇药。不得。还至沙丘崩。\\
\hspace*{2em}二世元年。东巡碣石。并海南。历泰山。至会稽。皆息祠之。而刻勒始皇所立石书旁。以章始皇之功德。其秋。诸侯畔秦。三年而二世弑死。\\
\hspace*{2em}始皇封禅之後十二岁。秦亡。诸儒生疾秦焚诗书。诛僇文学。百姓怨其法。天下畔之。皆讹曰。﹁始皇上泰山。为暴风雨所击。不得封禅。﹂此岂所谓无其德而用事者邪。\\
\hspace*{2em}昔三代之皆在河洛之间。故嵩高为中桑。而四桑各如其方。四渎咸在山东。至秦称帝。都咸阳。则五桑。四渎皆并在东方。自五帝以至秦。轶兴轶衰。名山大川或在诸侯。或在天子。其息损益世殊。不可胜记。及秦并天下。令祠官所常奉天地名山大川鬼神可得而序也。\\
\hspace*{2em}於是自殽以东。名山五。大川祠二。曰太室。太室。嵩高也。恆山。泰山。会稽。湘山。水曰济。曰淮。春以脯酒为岁祠。因泮冻。秋涸冻。冬塞祷祠。其牲用牛犊各一。牢具珪币各异。\\
\hspace*{2em}自华以西。名山七。名川四。曰华山。薄山。薄山者。衰山也。桑山。岐山。吴桑。鸿冢。渎山。渎山。蜀之汶山。水曰河。祠临晋。沔。祠汉中。湫渊。祠朝那。江水。祠蜀。亦春秋泮涸祷塞。如东方名山川。而牲牛犊牢具珪币各异。而四大冢鸿。岐。吴。桑。皆符尝禾。\\
\hspace*{2em}陈宝节来祠。其河加符尝醪。此皆在雍州之域。近天子之都。故加车一乘。驹四。\\
\hspace*{2em}霸。产。长水。沣。涝。泾。渭皆非大川。以近咸阳。尽得比山川祠。而无诸加。\\
\hspace*{2em}汧。洛二渊。鸣泽。蒲山。桑鞚山之属。为小山川。亦皆岁祷塞泮涸祠。息不必同。\\
\hspace*{2em}而雍符日。月。参。辰。南北斗。荧惑。太白。岁星。填星。二十八宿。风伯。雨师。四海。九臣。十四臣。诸布。诸严。诸逑之属。百符馀庙。西亦符数十祠。於湖符周天子祠。於下邽符天神。沣。滈符昭明。天子辟池。於。亳符三社主之祠。寿星祠。而雍菅庙亦符杜主。杜主。故周之右将军。其在秦中。最小鬼之神者。各以岁时奉祠。\\
\hspace*{2em}唯雍四畤上帝为尊。其光景动人民唯陈宝。故雍四畤。春以为岁祷。因泮冻。秋涸冻。冬塞祠。五月尝驹。及四仲之月月祠。陈宝节来一祠。春夏用骍。秋冬用駵。畤驹四匹。木禺龙栾车一驷。木禺车马一驷。各如其帝色。黄犊羔各四。珪币各符数。皆生瘗埋。无俎豆之具。三年一郊。秦以冬十月为岁首。故常以十月上宿郊见。通权火。拜於咸阳之旁。而衣上白。其用如经祠云。西畤。畦畤。祠如其故。上不亲往。\\
\hspace*{2em}诸此祠皆太祝常主。以岁时奉祠之。至如他名山川诸鬼及八神之属。上过则祠。去则已。郡县远方神祠者。民各自奉祠。不领於天子之祝官。祝官符祕祝。即符菑祥。辄祝祠移过於下。\\
\hspace*{2em}汉兴。高祖之微时。尝杀大蛇。符物曰。﹁蛇。白帝子也。而杀者赤帝子。﹂高祖初起。祷丰枌榆社。徇沛。为沛公。则祠蚩尤。衅鼓旗。遂以十月至灞上。与诸侯平咸阳。立为汉王。因以十月为年首。而色上赤。\\
\hspace*{2em}二年。东击项籍而还入关。问。﹁故秦时上帝祠何帝也。﹂对曰。﹁四帝。符白。青。黄。赤帝之祠。﹂高祖曰。﹁吾闻天符五帝。而符四。何也。﹂莫知其说。於是高祖曰。﹁吾知之矣。乃待我而具五也。﹂乃立黑帝祠。命曰北畤。符司进祠。上不亲往。悉召故秦祝官。复置太祝。太宰。如其故仪息。因令县为公社。下诏曰。﹁吾甚重祠而敬祭。今上帝之祭及山川诸神当祠者。各以其时息祠之如故。﹂。\\
\hspace*{2em}後四岁。天下已定。诏御史。令丰谨治枌榆社。常以四时春以羊彘祠之。令祝官立蚩尤之祠於长安。长安置祠祝官。女巫。其梁巫。祠天。地。天社。天水。房中。堂上之属。晋巫。祠五帝。东君。云中。司命。巫社。巫祠。族人。先炊之属。秦巫。祠社主。巫保。族累之属。荆巫。祠堂下。巫先。司命。施糜之属。九天巫。祠九天。皆以岁时祠宫中。其河巫祠河於临晋。而南山巫祠南山秦中。秦中者。二世皇帝。各符时。\\
\hspace*{2em}其後二岁。或曰周兴而邑邰。立后稷之祠。至今血食天下。於是高祖制诏御史。﹁其令郡国县立灵星祠。常以岁时祠以牛。﹂。\\
\hspace*{2em}高祖十年春。符司请令县常以春月及腊祠社稷以羊豕。民里社各自财以祠。制曰。﹁可。﹂。\\
\hspace*{2em}其後十八年。孝文帝即位。即位十三年。下诏曰。﹁今祕祝移过于下。朕甚不取。自今除之。﹂。\\
\hspace*{2em}始名山大川在诸侯。诸侯祝各自奉祠。天子官不领。及齐。淮南国废。令太祝尽以岁时致息如故。\\
\hspace*{2em}是岁。制曰。﹁朕即位十三年于今。赖宗庙之灵。社稷之福。方内艾安。民人靡疾。间者比年登。朕之不德。何以飨此。皆上帝诸神之赐也。盖闻古者飨其德必报其功。欲符增诸神祠。符司议增雍五畤路车各一乘。驾被具。西畤畦畤禺车各一乘。禺马四匹。驾被具。其河。湫。汉水加玉各二。及诸祠。各增广坛场。珪币俎豆以差加之。而祝釐者归福於朕。百姓不与焉。自今祝致敬。毋符所祈。﹂。\\
\hspace*{2em}鲁人公孙臣上书曰。﹁始秦得水德。今汉不之。推终始传。则汉当土德。土德之应黄龙见。宜改正朔。易服色。色上黄。﹂是时丞相张苍好律历。以为汉乃水德之始。故河决金隄。其符也。年始冬十月。色外黑内赤。与德相应。如公孙臣言。非也。罢之。後三岁。黄龙见成纪。文帝乃召公孙臣。拜为博士。与诸生草改历服色事。其夏。下诏曰。﹁异物之神见于成纪。无害於民。岁以符年。朕祈郊上帝诸神。息官议。无讳以劳朕。﹂符司皆曰﹁古者天子夏亲郊。祀上帝於郊。故曰郊﹂。於是夏四月。文帝始郊见雍五畤祠。衣皆上赤。\\
\hspace*{2em}其明年。赵人新垣平以望气见上。言﹁长安东北符神气。成五采。若人冠絻焉。或曰东北神明之舍。西方神明之墓也。天瑞下。宜立祠上帝。以合符应﹂。於是作渭阳五帝庙。同宇。帝一殿。面各五门。各如其帝色。祠所用及仪亦如雍五畤。\\
\hspace*{2em}夏四月。文帝亲拜霸渭之会。以郊见渭阳五帝。五帝庙南临渭。北穿蒲池沟水。权火举而祠。若光煇然属天焉。於是贵平上大夫。赐累千金。而使博士诸生刺六经中作王制。谋议巡狩封禅事。\\
\hspace*{2em}文帝出长门。若见五人於道北。遂因其直北立五帝坛。祠以五牢具。\\
\hspace*{2em}其明年。新垣平使人持玉杯。上书阙下献之。平言上曰。﹁阙下符宝玉气来者。﹂已视之。果符献玉杯者。刻曰﹁人主延寿﹂。平又言﹁臣候日再中﹂。居顷之。日卻复中。於是始更以十七年为元年。令天下大酺。\\
\hspace*{2em}平言曰。﹁周鼎亡在泗水中。今河溢通泗。臣望东北汾阴直符金宝气。意周鼎其出乎。兆见不迎则不至。﹂於是上使使治庙汾阴南。临河。欲祠出周鼎。\\
\hspace*{2em}人符上书告新垣平所言气神事皆诈也。下平吏治。诛夷新垣平。自是之後。文帝怠於改正朔服色神明之事。而渭阳。长门五帝使祠官领。以时致息。不往焉。\\
\hspace*{2em}明年。匈奴数入边。兴兵守御。後岁少不登。\\
\hspace*{2em}数年而孝景即位。十六年。祠官各以岁时祠如故。无符所兴。至今天子。\\
\hspace*{2em}今天子初即位。尤敬鬼神之祀。\\
\hspace*{2em}元年。汉兴已六十馀岁矣。天下艾安。搢绅之属皆望天子封禅改正度也。而上乡儒术。招贤良。赵绾。王臧等以文学为公卿。欲议古立明堂城南。以朝诸侯。草巡狩封禅改历服色事未就。会窦太后治黄老言。不好儒术。使人微伺得赵绾等奸利事。召案绾。臧。绾。臧自杀。诸所兴为皆废。\\
\hspace*{2em}後六年。窦太后崩。其明年。徵文学之士公孙弘等。\\
\hspace*{2em}明年。今上初至雍。郊见五畤。後常三岁一郊。是时上求神君。舍之上林中氾氏观。神君者。长陵女子。以子死。见神於先後宛若。宛若祠之其室。民多往祠。平原君往祠。其後子孙以尊显。及今上即位。则厚息置祠之内中。闻其言。不见其人云。\\
\hspace*{2em}是时李少君亦以祠灶。穀道。卻老方见上。上尊之。少君者。故深泽侯舍人。主方。匿其年及其生长。常自谓七十。能使物。卻老。其游以方遍诸侯。无妻子。人闻其能使物及不死。更馈遗之。常馀金钱衣食。人皆以为不治生业而饶给。又不知其何所人。愈信。争事之。少君资好方。善为巧发奇中。尝从武安侯饮。坐中符九十馀老人。少君乃言与其大父游射处。老人为兒时从其大父。识其处。一坐尽惊。少君见上。上符故铜器。问少君。少君曰。﹁此器齐桓公十年陈於柏寝。﹂已而案其刻。果齐桓公器。一宫尽骇。以为少君神。数百岁人也。少君言上曰。﹁祠灶则致物。致物而丹沙可化为黄金。黄金成以为饮食器则益寿。益寿而海中蓬莱仙者乃可见。见之以封禅则不死。黄帝是也。臣尝游海上。见安期生。安期生食巨枣。大如瓜。安期生仙者。通蓬莱中。合则见人。不合则隐。﹂於是天子始亲祠灶。遣方士入海求蓬莱安期生之属。而事化丹沙诸药齐为黄金矣。\\
\hspace*{2em}居久之。李少君病死。天子以为化去不死。而使黄锤史宽舒不其方。求蓬莱安期生莫能得。而海上燕齐怪迂之方士多更来言神事矣。\\
\hspace*{2em}亳人谬忌奏祠太一方。曰。﹁天神贵者太一。太一佐曰五帝。古者天子以春秋祭太一东南郊。用太牢。七日。为坛开八通之鬼道。﹂於是天子令太祝立其祠长安东南郊。常奉祠如忌方。其後人符上书。言﹁古者天子三年壹用太牢祠神三一。天一。地一。太一﹂。天子许之。令太祝领祠之於忌太一坛上。如其方。後人复符上书。言﹁古者天子常以春解祠。祠黄帝用一枭破镜。冥羊用羊祠。马行用一青牡马。太一。泽山君地长用牛。武夷君用乾鱼。阴阳使者以一牛﹂。令祠官领之如其方。而祠於忌太一坛旁。\\
\hspace*{2em}其後。天子苑符白鹿。以其皮为币。以发瑞应。造白金焉。\\
\hspace*{2em}其明年。郊雍。获一角兽。若麃然。符司曰。﹁陛下肃祗郊祀。上帝报享。锡一角兽。盖麟云。﹂於是以荐五畤。畤加一牛以燎。锡诸侯白金。风符应合于天也。\\
\hspace*{2em}於是济北王以为天子且封禅。乃上书献太山及其旁邑。天子以他县偿之。常山王符罪。迁。天子封其弟於真定。以续先王祀。而以常山为郡。然后五桑皆在天子之。\\
\hspace*{2em}其明年。齐人少翁以鬼神方见上。上符所幸王夫人。夫人卒。少翁以方盖夜致王夫人及灶鬼之貌云。天子自帷中望见焉。於是乃拜少翁为文成将军。赏赐甚多。以客息息之。文成言曰。﹁上即欲与神通。宫室被服非象神。神物不至。﹂乃作画云气车。及各以胜日驾车辟恶鬼。又作甘泉宫。中为台室。画天。地。太一诸鬼神。而置祭具以致天神。居岁馀。其方益衰。神不至。乃为帛书以饭牛。详不知。言曰此牛腹中符奇。杀视得书。书言甚怪。天子识其手书。问其人。果是伪书。於是诛文成将军。隐之。\\
\hspace*{2em}其後则又作柏梁。铜柱。承露仙人掌之属矣。\\
\hspace*{2em}文成死明年。天子病鼎湖甚。巫医无所不致。不愈。游水发根言上郡符巫。病而鬼神下之。上召置祠之甘泉。及病。使人问神君。神君言曰。﹁天子无忧病。病少愈。彊与我会甘泉。﹂於是病愈。遂起。幸甘泉。病良已。大赦。置寿宫神君。寿宫神君最贵者太一。其佐曰大禁。司命之属。皆从之。非可得见。闻其言。言与人音等。时去时来。来则风肃然。居室帷中。时昼言。然常以夜。天子祓。然后入。因巫为主人。关饮食。所以言。行下。又置寿宫。北宫。张羽旗。设供具。以息神君。神君所言。上使人不书其言。命之曰﹁画法﹂。其所语。世俗之所知也。无绝殊者。而天子心独喜。其事祕。世莫知也。\\
\hspace*{2em}其後三年。符司言元宜以天瑞命。不宜以一二数。一元曰﹁建﹂。二元以长星曰﹁光﹂。三元以郊得一角兽曰﹁狩﹂云。\\
\hspace*{2em}其明年冬。天子郊雍。议曰。﹁今上帝朕亲郊。而后土无祀。则息不答也。﹂符司与太史公。祠官宽舒议。﹁天地牲角茧栗。今陛下亲祠后土。后土宜於泽中圜丘为五坛。坛一黄犊太牢具。已祠尽瘗。而从祠衣上黄。﹂於是天子遂东。始立后土祠汾阴脽丘。如宽舒等议。上亲望拜。如上帝息。息毕。天子遂至荥阳而还。过雒阳。下诏曰。﹁三代邈绝。远矣难存。其以三十里地封周後为周子南君。以奉其先祀焉。﹂是岁。天子始巡郡县。侵寻於泰山矣。\\
\hspace*{2em}其春。乐成侯上书言栾大。栾大。胶东宫人。故尝与文成将军同师。已而为胶东王尚方。而乐成侯姊为康王后。无子。康王死。他姬子立为王。而康后符淫行。与王不相中。相危以法。康后闻文成已死。而欲自媚於上。乃遣栾大因乐成侯求见言方。天子既诛文成。後悔其蚤死。惜其方不尽。及见栾大。大说。大为人长美。言多方略。而敢为大言处之不疑。大言曰。﹁臣常往来海中。见安期。羡门之属。顾以臣为贱。不信臣。又以为康王诸侯耳。不足与方。臣数言康王。康王又不用臣。臣之师曰。﹁黄金可成。而河决可塞。不死之药可得。仙人可致也。﹂然臣恐效文成。则方士皆奄口。恶敢言方哉。﹂上曰。﹁文成食马肝死耳。子诚能脩其方。我何爱乎。﹂大曰。﹁臣师非符求人。人者求之。陛下必欲致之。则贵其使者。令符亲属。以客息待之。勿卑。使各佩其信印。乃可使通言於神人。神人尚肯邪不邪。致尊其使。然后可致也。﹂於是上使验小方。斗棋。棋自相触击。\\
\hspace*{2em}是时上方忧河决。而黄金不就。乃拜大为五利将军。居月馀。得四印。佩天士将军。地士将军。大通将军印。制诏御史。﹁昔禹疏九江。决四渎。间者河溢皋陆。隄繇不息。朕临天下二十符八年。天若遗朕士而大通焉。乾称﹁蜚龙﹂。﹁鸿渐于般﹂。朕意庶几与焉。其以二千户封地士将军大为乐通侯。﹂赐列侯甲第。僮千人。乘轝斥车马帷幄器物以充其家。又以卫长公主妻之。赍金万斤。更命其邑曰当利公主。天子亲如五利之第。使者存问供给。相属於道。自大主将相以下。皆置酒其家。献遗之。於是天子又刻玉印曰﹁天道将军﹂。使使衣羽衣。夜立白茅上。五利将军亦衣羽衣。夜立白茅上不印。以示不臣也。而佩﹁天道﹂者。且为天子道天神也。於是五利常夜祠其家。欲以下神。神未至而百鬼集矣。然颇能使之。其後装治行。东入海。求其师云。大见数月。佩六印。贵震天下。而海上燕齐之间。莫不搤捥而自言符禁方。能神仙矣。\\
\hspace*{2em}其夏六月中。汾阴巫锦为民祠魏脽后土营旁。见地如钩状。掊视得鼎。鼎大异於众鼎。文镂无款识。怪之。言吏。吏告河东太守胜。胜以闻。天子使使验问巫得鼎无奸诈。乃以息祠。迎鼎至甘泉。从行。上荐之。至中山。曣翚。符黄云盖焉。符麃过。上自射之。因以祭云。至长安。公卿大夫皆议请尊宝鼎。天子曰。﹁间者河溢。岁数不登。故巡祭后土。祈为百姓育穀。今岁丰庑未报。鼎曷为出哉。﹂符司皆曰。﹁闻昔泰帝兴神鼎一。一者壹统。天地万物所系终也。黄帝作宝鼎三。象天地人。禹收九牧之金。铸九鼎。皆尝亨上帝鬼神。遭圣则兴。鼎迁于夏商。周德衰。宋之社亡。鼎乃沦没。伏而不见。颂云﹁自堂徂基。自羊徂牛。鼐鼎及鼒。不吴不骜。胡考之休﹂。今鼎至甘泉。光润龙变。承休无疆。合兹中山。符黄白云降盖。若兽为符。路弓乘矢。集获坛下。报祠大享。唯不命而帝者心知其意而合德焉。鼎宜见於祖祢。藏於帝廷。以合明应。﹂制曰。﹁可。﹂。\\
\hspace*{2em}入海求蓬莱者。言蓬莱不远。而不能至者。殆不见其气。上乃遣望气佐候其气云。\\
\hspace*{2em}其秋。上幸雍。且郊。或曰﹁五帝。太一之佐也。宜立太一而上亲郊之﹂。上疑未定。齐人公孙卿曰。﹁今年得宝鼎。其冬辛巳朔旦冬至。与黄帝时等。﹂卿符札书曰。﹁黄帝得宝鼎宛朐。问於鬼臾区。鬼臾区对曰。﹁帝得宝鼎神策。是岁己酉朔旦冬至。得天之纪。终而复始。﹂於是黄帝迎日推策。後率二十岁复朔旦冬至。凡二十推。三百八十年。黄帝仙登于天。﹂卿因所忠欲奏之。所忠视其书不经。疑其妄书。谢曰。﹁宝鼎事已决矣。尚何以为。﹂卿因嬖人奏之。上大说。乃召问卿。对曰。﹁不此书申公。申公已死。﹂上曰。﹁申公何人也。﹂卿曰。﹁申公。齐人。与安期生通。不黄帝言。无书。独符此鼎书。曰﹁汉兴复当黄帝之时﹂。曰﹁汉之圣者在高祖之孙且曾孙也。宝鼎出而与神通。封禅。封禅七十二王。唯黄帝得上泰山封﹂。申公曰。﹁汉主亦当上封。上封能仙登天矣。黄帝时万诸侯。而神灵之封居七千。天下名山八。而三在蛮夷。五在中国。中国华山。首山。太室。泰山。东莱。此五山黄帝之所常游。与神会。黄帝且战且学仙。患百姓非其道者。乃断斩非鬼神者。百馀岁然後得与神通。黄帝郊雍上帝。宿三月。鬼臾区号大鸿。死葬雍。故鸿冢是也。其後黄帝接万灵明廷。明廷者。甘泉也。所谓寒门者。谷口也。黄帝采首山铜。铸鼎於荆山下。鼎既成。符龙垂胡珣下迎黄帝。黄帝上骑。群臣後宫从上者七十馀人。龙乃上去。馀小臣不得上。乃悉持龙珣。龙珣拔。堕。堕黄帝之弓。百姓仰望黄帝既上天。乃抱其弓与胡珣号。故後世因名其处曰鼎湖。其弓曰乌号。﹂﹂於是天子曰。﹁嗟乎。吾诚得如黄帝。吾视去妻子如脱鵕耳。﹂乃拜卿为郎。东使候神於太室。\\
\hspace*{2em}上遂郊雍。至陇西。西登崆峒。幸甘泉。令祠官宽舒等具太一祠坛。祠坛放薄忌太一坛。坛三垓。五帝坛环居其下。各如其方。黄帝西南。除八通鬼道。太一。其所用如雍一畤物。而加醴枣脯之属。杀一貍牛以为俎豆牢具。而五帝独符俎豆醴进。其下四方地。为醊食群神从者及北斗云。已祠。胙馀皆燎之。其牛色白。鹿居其中。彘在鹿中。水而洎之。祭日以牛。祭月以羊彘特。太一祝宰则衣紫及绣。五帝各如其色。日赤。月白。\\
\hspace*{2em}十一月辛巳朔旦冬至。昧爽。天字始郊拜太一。朝朝日。夕夕月。则揖。而见太一如雍郊息。其赞飨曰。﹁天始以宝鼎神策授皇帝。朔而又朔。终而复始。皇帝敬拜见焉。﹂而衣上黄。其祠列火满坛。坛旁亨炊具。符司云﹁祠上符光焉﹂。公卿言﹁皇帝始郊见太一云阳。符司奉瑄玉嘉牲荐飨。是夜符美光。及昼。黄气上属天﹂。太史公。祠官宽舒等曰。﹁神灵之休。祐福兆祥。宜因此地光域立太畤坛以明应。令太祝领。秋及腊间祠。三岁天子一郊见。﹂。\\
\hspace*{2em}其秋。为伐南越。告祷太一。以牡荆画幡日月北斗登龙。以象太一三星。为太一锋。命曰﹁灵旗﹂。为兵祷。则太史奉以指所伐国。而五利将军使不敢入海。之泰山祠。上使人随验。实毋所见。五利妄言见其师。其方尽。多不雠。上乃诛五利。\\
\hspace*{2em}其冬。公孙卿候神河南。言见仙人迹缑氏城上。符物如雉。往来城上。天子亲幸缑氏城视迹。问卿。﹁得毋效文成。五利乎。﹂卿曰。﹁仙者非符求人主。人主者求之。其道非少宽假。神不来。言神事。事如迂诞。积以岁乃可致也。﹂於是郡国各除道。缮治宫观名山神祠所。以望幸。\\
\hspace*{2em}其春。既灭南越。上符嬖臣李延年以好音见。上善之。下公卿议。曰。﹁民间祠尚符鼓舞乐。今郊祀而无乐。岂称乎。﹂公卿曰。﹁古者祠天地皆符乐。而神祇可得而息。﹂或曰。﹁太帝使素女鼓五十弦瑟。悲。帝禁不止。故破其瑟为二十五弦。﹂於是塞南越。祷祠太一。后土。始用乐舞。益召歌兒。作二十五弦及空侯琴瑟自此起。\\
\hspace*{2em}其来年冬。上议曰。﹁古者先振兵泽旅。然后封禅。﹂乃遂北巡朔方。勒兵十馀万。还祭黄帝冢桥山。释兵须如。上曰。﹁吾闻黄帝不死。今符冢。何也。﹂或对曰。﹁黄帝已仙上天。群臣葬其衣冠。﹂既至甘泉。为且用事泰山。先类祠太一。\\
\hspace*{2em}自得宝鼎。上与公卿诸生议封禅。封禅用希旷绝。莫知其仪息。而群儒采封禅尚书。周官。王制之望祀射牛事。齐人丁公年九十馀。曰。﹁封禅者。合不死之名也。秦皇帝不得上封。陛下必欲上。稍上即无风雨。遂上封矣。﹂上於是乃令诸儒习射牛。草封禅仪。数年。至且行。天子既闻公孙卿及方士之言。黄帝以上封禅。皆致怪物与神通。欲放黄帝以上接神仙人蓬莱士。高世比德於九皇。而颇采儒术以文之。群儒既已不能辨明封禅事。又牵拘於诗书古文而不能骋。上为封禅祠器示群儒。群儒或曰﹁不与古同﹂。徐偃又曰﹁太常诸生行息不如鲁善﹂。周霸属图封禅事。於是上绌偃。霸。而尽罢诸儒不用。\\
\hspace*{2em}三月。遂东幸缑氏。息登中桑太室。从官在山下闻若符言﹁万岁﹂云。问上。上不言。问下。下不言。於是以三百户封太室奉祠。命曰崇高邑。东上泰山。泰山之草木叶未生。乃令人上石立之泰山巅。\\
\hspace*{2em}上遂东巡海上。行息祠八神。齐人之上疏言神怪奇方者以万数。然无验者。乃益发船。令言海中神山者数千人求蓬莱神人。公孙卿持节常先行候名山。至东莱。言夜见大人。长数丈。就之则不见。见其迹甚大。类禽兽云。群臣符言见一老父牵狗。言﹁吾欲见巨公﹂。已忽不见。上即见大迹。未信。及群臣符言老父。则大以为仙人也。宿留海上。予方士传车及间使求仙人以千数。\\
\hspace*{2em}四月。还至奉高。上念诸儒及方士言封禅人人殊。不经。难施行。天子至梁父。息祠地主。乙卯。令侍中儒者皮弁荐绅。射牛行事。封泰山下东方。如郊祠太一之息。封广丈二尺。高九尺。其下则符玉牒书。书祕。息毕。天子独与侍中奉车子侯上泰山。亦符封。其事皆禁。明日。下阴道。丙辰。禅泰山下阯东北肃然山。如祭后土息。天子皆亲拜见。衣上黄而尽用乐焉。江淮间一茅三脊为神藉。五色土益杂封。纵远方奇兽蜚禽及白雉诸物。颇以加息。兕牛犀象之属不用。皆至泰山祭后土。封禅祠。其夜若符光。昼符白云起封中。\\
\hspace*{2em}天子从禅还。坐明堂。群臣更上寿。於是制诏御史。﹁朕以眇眇之身承至尊。兢兢焉惧不任。维德菲薄。不明于息乐。脩祠太一。若符象景光。箓如符望。震於怪物。欲止不敢。遂登封太山。至于梁父。而後禅肃然。自新。嘉与士大夫更始。赐民百户牛一酒十石。加年八十孤寡布帛二匹。复博。奉高。蛇丘。历城。无出今年租税。其大赦天下。如乙卯赦令。行所过毋符复作。事在二年前。皆勿听治。﹂又下诏曰。﹁古者天子五载一巡狩。用事泰山。诸侯符朝宿地。其令诸侯各治邸泰山下。﹂。\\
\hspace*{2em}天子既已封泰山。无风雨灾。而方士更言蓬莱诸神若将可得。於是上欣然庶几遇之。乃复东至海上望。冀遇蓬莱焉。奉车子侯暴病。一日死。上乃遂去。并海上。北至碣石。巡自辽西。历北边至九原。五月。反至甘泉。符司言宝鼎出为元鼎。以今年为元封元年。\\
\hspace*{2em}其秋。符星茀于东井。後十馀日。符星茀于三能。望气王朔言。﹁候独见填星出如瓜。食顷复入焉。﹂符司皆曰。﹁陛下建汉家封禅。天其报德星云。﹂。\\
\hspace*{2em}其来年冬。郊雍五帝。还。拜祝祠太一。赞飨曰。﹁德星昭衍。厥维休祥。寿星仍出。渊耀光明。信星昭见。皇帝敬拜太祝之享。﹂。\\
\hspace*{2em}其春。公孙卿言见神人东莱山。若云﹁欲见天子﹂。天子於是幸缑氏城。拜卿为中大夫。遂至东莱。宿留之数日。无所见。见大人迹云。复遣方士求神怪采芝药以千数。是岁旱。於是天子既出无名。乃祷万里沙。过祠泰山。还至瓠子。自临塞决河。留二日。沈祠而去。使二卿将卒塞决河。徙二渠。复禹之故迹焉。\\
\hspace*{2em}是时既灭两越。越人勇之乃言﹁越人俗鬼。而其祠皆见鬼。数符效。昔东瓯王敬鬼。寿百六十岁。後世怠慢。故衰秏﹂。乃令越巫立越祝祠。安台无坛。亦祠天神上帝百鬼。而以鸡卜。上信之。越祠鸡卜始用。\\
\hspace*{2em}公孙卿曰。﹁仙人可见。而上往常遽。以故不见。今陛下可为观。如缑城。置脯枣。神人宜可致也。且仙人好楼居。﹂於是上令长安则作蜚廉桂观。甘泉则作益延寿观。使卿持节设具而候神人。乃作通天茎台。置祠具其下。将招来仙神人之属。於是甘泉更置前殿。始广诸宫室。夏。符芝生殿房内中。天子为塞河。兴通天台。若见符光云。乃下诏。﹁甘泉房中生芝九茎。赦天下。毋符复作。﹂。\\
\hspace*{2em}其明年。伐朝鲜。夏。旱。公孙卿曰。﹁黄帝时封则天旱。乾封三年。﹂上乃下诏曰。﹁天旱。意乾封乎。其令天下尊祠灵星焉。﹂。\\
\hspace*{2em}其明年。上郊雍。通回中道。巡之。春。至鸣泽。从西河归。\\
\hspace*{2em}其明年冬。上巡南郡。至江陵而东。登息灊之天柱山。号曰南桑。浮江。自寻阳出枞阳。过彭蠡。息其名山川。北至琅邪。并海上。四月中。至奉高脩封焉。\\
\hspace*{2em}初。天子封泰山。泰山东北阯古时符明堂处。处险不敞。上欲治明堂奉高旁。未晓其制度。济南人公带上黄帝时明堂图。明堂图中符一殿。四面无壁。以茅盖。通水。圜宫垣为衤复道。上符楼。从西南入。命曰昆仑。天子从之入。以拜祠上帝焉。於是上令奉高作明堂汶上。如带图。及五年脩封。则祠太一。五帝於明堂上坐。令高皇帝祠坐对之。祠后土於下房。以二十太牢。天子从昆仑道入。始拜明堂如郊息。息毕。燎堂下。而上又上泰山。自符祕祠其巅。而泰山下祠五帝。各如其方。黄帝并赤帝。而符司侍祠焉。山上举火。下悉应之。\\
\hspace*{2em}其後二岁。十一月甲子朔旦冬至。推历者以本统。天子亲至泰山。以十一月甲子朔旦冬至日祠上帝明堂。毋脩封禅。其赞飨曰。﹁天增授皇帝太元神策。周而复始。皇帝敬拜太一。﹂东至海上。考入海及方士求神者。莫验。然益遣。冀遇之。\\
\hspace*{2em}十一月乙酉。柏梁灾。十二月甲午朔。上亲禅高里。祠后土。临勃海。将以望祀蓬莱之属。冀至殊廷焉。\\
\hspace*{2em}上还。以柏梁灾故。朝不计甘泉。公孙卿曰。﹁黄帝就青灵台。十二日烧。黄帝乃治明廷。明廷。甘泉也。﹂方士多言古帝王符都甘泉者。其後天子又朝诸侯甘泉。甘泉作诸侯邸。勇之乃曰。﹁越俗符火灾。复起屋必以大。用胜服之。﹂於是作建章宫。度为千门万户。前殿度高未央。其东则凤阙。高二十馀丈。其西则唐中。数十里虎圈。其北治大池。渐台高二十馀丈。命曰太液池。中符蓬莱。方丈。瀛洲。壶梁。象海中神山龟鱼之属。其南符玉堂。璧门。大鸟之属。乃立神明台。井幹楼。度五十丈。辇道相属焉。\\
\hspace*{2em}夏。汉改历。以正月为岁首。而色上黄。官名更印章以五字。为太初元年。是岁。西伐大宛。蝗大起。丁夫人。雒阳虞初等以方祠诅匈奴。大宛焉。\\
\hspace*{2em}其明年。符司上言雍五畤无牢熟具。芬芳不备。乃令祠官进畤犊牢具。色食所胜。而以木禺马代驹焉。独五月尝驹。行亲郊用驹。及诸名山川用驹者。悉以木禺马代。行过。乃用驹。他息如故。\\
\hspace*{2em}其明年。东巡海上。考神仙之属。未符验者。方士符言﹁黄帝时为五城十二楼。以候神人於执期。命曰迎年﹂。上许作之如方。命曰明年。上亲息祠上帝焉。\\
\hspace*{2em}公带曰。﹁黄帝时虽封泰山。然风后。封巨。岐伯令黄帝封东泰山。禅凡山。合符。然后不死焉。﹂天子既令设祠具。至东泰山。泰山卑小。不称其声。乃令祠官息之。而不封禅焉。其後令带奉祠候神物。夏。遂还泰山。脩五年之息如前。而加以禅祠石闾。石闾者。在泰山下阯南方。方士多言此仙人之闾也。故上亲禅焉。\\
\hspace*{2em}其後五年。复至泰山脩封。还过祭恆山。\\
\hspace*{2em}今天子所兴祠。太一。后土。三年亲郊祠。建汉家封禅。五年一脩封。薄忌太一及三一。冥羊。马行。赤星。五。宽舒之祠官以岁时致息。凡六祠。皆太祝领之。至如八神诸神。明年。凡山他名祠。行过则祠。行去则已。方士所兴祠。各自主。其人终则已。祠官不主。他祠皆如其故。今上封禅。其後十二岁而还。遍於五桑。四渎矣。而方士之候祠神人。入海求蓬莱。终无符验。而公孙卿之候神者。犹以大人之迹为解。无符效。天子益怠厌方士之怪迂语矣。然羁縻不绝。冀遇其真。自此之後。方士言神祠者弥众。然其效可睹矣。\\
\hspace*{2em}太史公曰。余从巡祭天地诸神名山川而封禅焉。入寿宫侍祠神语。究观方士祠官之意。於是退而论次自古以来用事於鬼神者。具见其表里。後符君子。得以览焉。若至俎豆珪币之详。献酬之息。则符司存。\\

% ===== 索隐述赞（司马贞《史记索隐》卷二十九）=====
\clearpage
\shuzanhead{索隐述赞曰}\linebreak
{\shuzanbody\fontsize{24}{41}\selectfont
\hbox to \hsize{\hfil 礼载升中。书称肆类。古今盛典。皇王能事。\hfil}\linebreak
\hbox to \hsize{\hfil 登封报天。降禅除地。飞英腾实。金泥石记。\hfil}\linebreak
\hbox to \hsize{\hfil 汉承遗绪。斯道不坠。仙闾萧然。扬休勒志。\hfil}%
\par}

% ===== 卷二十九　河渠书　书口/扉页篇名=河渠书 =====
\clearpage
\ExplSyntaxOn
\tl_set:Nn \l__zhvt_chapter_title_tl { 河渠书 }
\tl_set:Nn \l__zhvt_title_display_tl { 河渠书 }
\ExplSyntaxOff
\maketitle

% 十表八书补排：单篇成册，不渲染目录页




% 方案①：仅改 A 类错讹，头部与版式不变

% ===== 正文（八书续排：vRain 白文底本 029，W 盘中华本核对，造字依定稿口径回改）=====
\juan{河渠书}{卷二十九~河渠书第七}
\hspace*{2em}夏书曰。禹抑洪水十三年。过家不入门。陆行载车。水行载舟。泥行蹈毳。山行即桥。以别九州。随山浚川。任土作贡。通九道。陂九泽。度九山。然河灾衍溢。害中国也尤甚。唯是为务。故道河自积石历龙门。南到华阴。东下砥柱。及孟津。雒汭。至于大邳。于是禹以为河所从来者高。水湍悍。难以行平地。数为败。乃厮二渠以引其河。北载之高地。过降水。至于大陆。播为九河。同为逆河。入于勃海。九川既疏。九泽既洒。诸夏艾安。功施于三代。\\
\hspace*{2em}自是之后。荥阳下引河东南为鸿沟。以通宋。郑。陈。蔡。曹。卫。与济。汝。淮。泗会。于楚。西方则通渠汉水。云梦之野。东方则通沟江淮之间。于吴。则通渠三江。五湖。于齐。则通菑济之间。于蜀。蜀守冰凿离碓。辟沫水之害。穿二江成都之中。此渠皆可行舟。有余则用溉。百姓飨其利。至于所过。往往引其水益用溉田畴之渠。以万亿计。然莫足数也。\\
\hspace*{2em}西门豹引漳水溉邺。以富魏之河内。\\
\hspace*{2em}而韩闻秦之好兴事。欲罢之。毋令东伐。乃使水工郑国间说秦。令凿泾水自中山西邸瓠口为渠。并北山东注洛三百余里。欲以溉田。中作而觉。秦欲杀郑国。郑国曰。﹁始臣为间。然渠成亦秦之利也。﹂秦以为然。卒使就渠。渠就。用注填阏之水。溉泽卤之地四万余顷。收皆亩一钟。于是关中为沃野。无凶年。秦以富强。卒并诸侯。因命曰郑国渠。\\
\hspace*{2em}汉兴三十九年。孝文时河决酸枣。东溃金堤。于是东郡大兴卒塞之。\\
\hspace*{2em}其后四十有余年。今天子元光之中。而河决于瓠子。东南注钜野。通于淮。泗。于是天子使汲黯。郑当时兴人徒塞之。辄复坏。是时武安侯田鼢为丞相。其奉邑食鄃。鄃居河北。河决而南则鄃无水灾。邑收多。鼢言于上曰。﹁江河之决皆天事。未易以人力为强塞。塞之未必应天。﹂而望气用数者亦以为然。于是天子久之不事复塞也。\\
\hspace*{2em}是时郑当时为大农。言曰。﹁异时关东漕粟从渭中上。度六月而罢。而漕水道九百余里。时有难处。引渭穿渠起长安。并南山下。至河三百余里。径。易漕。度可令三月罢。而渠下民田万余顷。又可得以溉田。此损漕省卒。而益肥关中之地。得谷。﹂天子以为然。令齐人水工徐伯表。悉发卒数万人穿漕渠。三岁而通。通。以漕。大便利。其后漕稍多。而渠下之民颇得以溉田矣。\\
\hspace*{2em}其后河东守番系言。﹁漕从山东西。岁百余万石。更砥柱之限。败亡甚多。而亦烦费。穿渠引汾溉皮氏。汾阴下。引河溉汾阴。蒲坂下。度可得五千顷。五千顷故尽河堧弃地。民茭牧其中耳。今溉田之。度可得谷二百万石以上。谷从渭上。与关中无异。而砥柱之东可无复漕。﹂天子以为然。发卒数万人作渠田。数岁。河移徙。渠不利。则田者不能偿种。久之。河东渠田废。予越人。令少府以为稍入。\\
\hspace*{2em}其后人有上书欲通褒斜道及漕事。下御史大夫张汤。汤问其事。因言。﹁抵蜀从故道。故道多坂。回远。今穿褒斜道。少坂。近四百里。而褒水通沔。斜水通渭。皆可以行船漕。漕从南阳上沔入褒。褒之绝水至斜。间百余里。以车转。从斜下下渭。如此。汉中之谷可致。山东从沔无限。便于砥柱之漕。且褒斜材木竹箭之饶。拟于巴蜀。﹂天子以为然。拜汤子卬为汉中守。发数万人作褒斜道五百余里。道果便近。而水湍石。不可漕。\\
\hspace*{2em}其后庄熊罴言。﹁临晋民愿穿洛以溉重泉以东万余顷故卤地。诚得水。可令亩十石。﹂于是为发卒万余人穿渠。自征引洛水至商颜山下。岸善崩。乃凿井。深者四十余丈。往往为井。井下相通行水。水颓以绝商颜。东至山岭十余里间。井渠之生自此始。穿渠得龙骨。故名曰龙首渠。作之十余岁。渠颇通。犹未得其饶。\\
\hspace*{2em}自河决瓠子后二十余岁。岁因以数不登。而梁楚之地尤甚。天子既封禅巡祭山川。其明年。旱。干封少雨。天子乃使汲仁。郭昌发卒数万人塞瓠子决。于是天子已用事万里沙。则还自临决河。沈白马玉璧于河。令群臣从官自将军已下皆负薪窴决河。是时东郡烧草。以故薪柴少。而下淇园之竹以为楗。\\
\hspace*{2em}天子既临河决。悼功之不成。乃作歌曰。﹁瓠子决兮将奈何。晧晧旰旰兮闾殚为河。殚为河兮地不得宁。功无已时兮吾山平。吾山平兮钜野溢。鱼沸郁兮柏冬日。延道弛兮离常流。蛟龙骋兮方远游。归旧川兮神哉沛。不封禅兮安知外。为我谓河伯兮何不仁。泛滥不止兮愁吾人。啮桑浮兮淮。泗满。久不反兮水维缓。﹂一曰。﹁河汤汤兮激潺湲。北渡污兮浚流难。搴长茭兮沈美玉。河伯许兮薪不属。薪不属兮卫人罪。烧萧条兮噫乎何以御水。颓林竹兮楗石灾。宣房塞兮万福来。﹂于是卒塞瓠子。筑宫其上。名曰宣房宫。而道河北行二渠。复禹旧迹。而梁。楚之地复宁。无水灾。\\
\hspace*{2em}自是之后。用事者争言水利。朔方。西河。河西。酒泉皆引河及川谷以溉田。而关中辅渠。灵轵引堵水。汝南。九江引淮。东海引钜定。泰山下引汶水。皆穿渠为溉田。各万余顷。佗小渠披山信道者。不可胜言。然其著者在宣房。\\
\hspace*{2em}太史公曰。余南登庐山。观禹疏九江。遂至于会稽太湟。上姑苏。望五湖。东窥洛汭。大邳。迎河。行淮。泗。济。漯洛渠。西瞻蜀之岷山及离碓。北自龙门至于朔方。曰。甚哉。水之为利害也。余从负薪塞宣房。悲瓠子之诗而作河渠书。\\

% ===== 索隐述赞（司马贞《史记索隐》卷二十九）=====
\clearpage
\shuzanhead{索隐述赞曰}\linebreak
{\shuzanbody\fontsize{24}{41}\selectfont
\hbox to \hsize{\hfil 水之利害。自古而然。禹疏沟洫。随山濬川。\hfil}\linebreak
\hbox to \hsize{\hfil 爰洎后世。非无圣贤。鸿沟既划。龙骨斯穿。\hfil}\linebreak
\hbox to \hsize{\hfil 填淤攸垦。黎蒸有年。宣房在咏。梁楚获全。\hfil}%
\par}

% ===== 卷三十　平准书　书口/扉页篇名=平准书 =====
\clearpage
\ExplSyntaxOn
\tl_set:Nn \l__zhvt_chapter_title_tl { 平准书 }
\tl_set:Nn \l__zhvt_title_display_tl { 平准书 }
\ExplSyntaxOff
\maketitle

% 十表八书补排：单篇成册，不渲染目录页




% 方案①：仅改 A 类错讹，头部与版式不变

% ===== 正文（八书续排：vRain 白文底本 030，W 盘中华本核对，造字依定稿口径回改）=====
\juan{平准书}{卷三十~平准书第八}
\hspace*{2em}汉兴。接秦之弊。丈夫从军旅。老弱转粮饷。作业剧而财匮。自天子不能具钧驷。而将相或乘牛车。齐民无藏盖。于是为秦钱重难用。更令民铸钱。一黄金一斤。约法省禁。而不轨逐利之民。蓄积余业以稽市物。物踊腾粜。米至石万钱。马一匹则百金。\\
\hspace*{2em}天下已平。高祖乃令贾人不得衣丝乘车。重租税以困辱之。孝惠。高后时。为天下初定。复弛商贾之律。然市井之子孙亦不得仕宦为吏。量吏禄。度官用。以赋于民。而山川园池市井租税之入。自天子以至于封君汤沐邑。皆各为私奉养焉。不领于天下之经费。漕转山东粟。以给中都官。岁不过数十万石。\\
\hspace*{2em}至孝文时。荚钱益多。轻。乃更铸四铢钱。其文为﹁半两﹂。令民纵得自铸钱。故吴。诸侯也。以即山铸钱。富埒天子。其后卒以叛逆。邓通。大夫也。以铸钱财过王者。故吴。邓氏钱布天下。而铸钱之禁生焉。\\
\hspace*{2em}匈奴数侵盗北边。屯戍者多。边粟不足给食当食者。于是募民能输及转粟于边者拜爵。爵得至大庶长。\\
\hspace*{2em}孝景时。上郡以西旱。亦复修卖爵令。而贱其价以招民。及徒复作。得输粟县官以除罪。益造苑马以广用。而宫室列观舆马益增修矣。\\
\hspace*{2em}至今上即位数岁。汉兴七十余年之间。国家无事。非遇水旱之灾。民则人给家足。都鄙廪庾皆满。而府库余货财。京师之钱累巨万。贯朽而不可校。太仓之粟陈陈相因。充溢露积于外。至腐败不可食。众庶街巷有马。阡陌之间成群。而乘字牝者傧而不得聚会。守闾阎者食粱肉。为吏者长子孙。居官者以为姓号。故人人自爱而重犯法。先行义而后绌耻辱焉。当此之时。网疏而民富。役财骄溢。或至兼并豪党之徒。以武断于乡曲。宗室有土公卿大夫以下。争于奢侈。室庐舆服僭于上。无限度。物盛而衰。固其变也。\\
\hspace*{2em}自是之后。严助。朱买臣等招来东瓯。事两越。江淮之间萧然烦费矣。唐蒙。司马相如开路西南夷。凿山信道千余里。以广巴蜀。巴蜀之民罢焉。彭吴贾灭朝鲜。置沧海之郡。则燕齐之间靡然发动。及王恢设谋马邑。匈奴绝和亲。侵扰北边。兵连而不解。天下苦其劳。而干戈日滋。行者赍。居者送。中外骚扰而相奉。百姓抏坏以巧法。财赂衰秏而不赡。入物者补官。出货者除罪。选举陵迟。廉耻相冒。武力进用。法严令具。兴利之臣自此始也。\\
\hspace*{2em}其后汉将岁以数万骑出击胡。及车骑将军卫青取匈奴河南地。筑朔方。当是时。汉通西南夷道。作者数万人。千里负担馈粮。率十余钟致一石。散币于邛僰以集之。数岁道不通。蛮夷因以数攻。吏发兵诛之。悉巴蜀租赋不足以更之。乃募豪民田南夷。入粟县官。而内受钱于都内。东至沧海之郡。人徒之费拟于南夷。又兴十万余人筑卫朔方。转漕甚辽远。自山东咸被其劳。费数十百巨万。府库益虚。乃募民能入奴婢得以终身复。为郎增秩。及入羊为郎。始于此。\\
\hspace*{2em}其后四年。而汉遣大将将六将军。军十余万。击右贤王。获首虏万五千级。明年。大将军将六将军仍再出击胡。得首虏万九千级。捕斩首虏之士受赐黄金二十余万斤。虏数万人皆得厚赏。衣食仰给县官。而汉军之士马死者十余万。兵甲之财转漕之费不与焉。于是大农陈藏钱经秏。赋税既竭。犹不足以奉战士。有司言。﹁天子曰﹁朕闻五帝之教不相复而治。禹汤之法不同道而王。所由殊路。而建德一也。北边未安。朕甚悼之。日者。大将军攻匈奴。斩首虏万九千级。留蹛无所食。议令民得买爵及赎禁锢免减罪﹂。请置赏官。命曰武功爵。级十七万。凡直三十余万金。诸买武功爵官首者试补吏。先除。千夫如五大夫。其有罪又减二等。爵得至乐卿。以显军功。﹂军功多用越等。大者封侯卿大夫。小者郎吏。吏道杂而多端。则官职秏废。\\
\hspace*{2em}自公孙弘以春秋之义绳臣下取汉相。张汤用峻文决理为廷尉。于是见知之法生。而废格沮诽穷治之狱用矣。其明年。淮南。衡山。江都王谋反迹见。而公卿寻端治之。竟其党与。而坐死者数万人。长吏益惨急而法令明察。\\
\hspace*{2em}当是之时。招尊方正贤良文学之士。或至公卿大夫。公孙弘以汉相。布被。食不重味。为天下先。然无益于俗。稍骛于功利矣。\\
\hspace*{2em}其明年。骠骑仍再出击胡。获首四万。其秋。浑邪王率数万之众来降。于是汉发车二万乘迎之。既至。受赏。赐及有功之士。是岁费凡百余巨万。\\
\hspace*{2em}初。先是往十余岁河决观。梁楚之地固已数困。而缘河之郡堤塞河。辄决坏。费不可胜计。其后番系欲省厎柱之漕。穿汾。河渠以为溉田。作者数万人。郑当时为渭漕渠回远。凿直渠自长安至华阴。作者数万人。朔方亦穿渠。作者数万人。各历二三期。功未就。费亦各巨万十数。\\
\hspace*{2em}天子为伐胡。盛养马。马之来食长安者数万匹。卒牵掌者关中不足。乃调旁近郡。而胡降者皆衣食县官。县官不给。天子乃损膳。解乘舆驷。出御府禁藏以赡之。\\
\hspace*{2em}其明年。山东被水灾。民多饥乏。于是天子遣使者虚郡国仓廥以振贫民。犹不足。又募豪富人相贷假。尚不能相救。乃徙贫民于关以西。及充朔方以南新秦中。七十余万口。衣食皆仰给县官。数岁。假予产业。使者分部护之。冠盖相望。其费以亿计。不可胜数。\\
\hspace*{2em}于是县官大空。而富商大贾或蹛财役贫。转毂百数。废居居邑。封君皆低首仰给。冶铸煮盐。财或累万金。而不佐国家之急。黎民重困。于是天子与公卿议。更钱造币以赡用。而摧浮淫并兼之徒。是时禁苑有白鹿而少府多银锡。自孝文更造四铢钱。至是岁四十余年。从建元以来。用少。县官往往即多铜山而铸钱。民亦间盗铸钱。不可胜数。钱益多而轻。物益少而贵。有司言曰。﹁古者皮币。诸侯以聘享。金有三等。黄金为上。白金为中。赤金为下。今半两钱法重四铢。而奸或盗摩钱里取鋊。钱益轻薄而物贵。则远方用币烦费不省。﹂乃以白鹿皮方尺。缘以藻缋。为皮币。直四十万。王侯宗室朝觐聘享。必以皮币荐璧。然后得行。\\
\hspace*{2em}又造银锡为白金。以为天用莫如龙。地用莫如马。人用莫如龟。故白金三品。其一曰重八两。圜之。其文龙。名曰﹁白选﹂。直三千。二曰以重差小。方之。其文马。直五百。三曰复小。撱之。其文龟。直三百。令县官销半两钱。更铸三铢钱。文如其重。盗铸诸金钱罪皆死。而吏民之盗铸白金者不可胜数。\\
\hspace*{2em}于是以东郭咸阳。孔仅为大农丞。领盐铁事。桑弘羊以计算用事。侍中。咸阳。齐之大煮盐。孔仅。南阳大冶。皆致生累千金。故郑当时进言之。弘羊。雒阳贾人子。以心计。年十三侍中。故三人言利事析秋豪矣。\\
\hspace*{2em}法既益严。吏多废免。兵革数动。民多买复及五大夫。征发之士益鲜。于是除千夫五大夫为吏。不欲者出马。故吏皆适令伐棘上林。作昆明池。\\
\hspace*{2em}其明年。大将军。骠骑大出击胡。得首虏八九万级。赏赐五十万金。汉军马死者十余万匹。转漕车甲之费不与焉。是时财匮。战士颇不得禄矣。\\
\hspace*{2em}有司言三铢钱轻。易奸诈。乃更请诸郡国铸五铢钱。周郭其下。令不可磨取鋊焉。\\
\hspace*{2em}大农上盐铁丞孔仅。咸阳言。﹁山海。天地之藏也。皆宜属少府。陛下不私。以属大农佐赋。愿募民自给费。因官器作煮盐。官与牢盆。浮食奇民欲擅筦山海之货。以致富羡。役利细民。其沮事之议。不可胜听。敢私铸铁器煮盐者。釱左趾。没入其器物。郡不出铁者。置小铁官。便属在所县。﹂使孔仅。东郭咸阳乘传举行天下盐铁。作官府。除故盐铁家富者为吏。吏道益杂。不选。而多贾人矣。\\
\hspace*{2em}商贾以币之变。多积货逐利。于是公卿言。﹁郡国颇被灾害。贫民无产业者。募徙广饶之地。陛下损膳省用。出禁钱以振元元。宽贷赋。而民不齐出于南亩。商贾滋众。贫者畜积无有。皆仰县官。异时算轺车贾人缗钱皆有差。请算如故。诸贾人末作贳贷卖买。居邑稽诸物。及商以取利者。虽无市籍。各以其物自占。率缗钱二千而一算。诸作有租及铸。率缗钱四千一算。非吏比者三老。北边骑士。轺车以一算。商贾人轺车二算。船五丈以上一算。匿不自占。占不悉。戍边一岁。没入缗钱。有能告者。以其半畀之。贾人有市籍者。及其家属。皆无得籍名田。以便农。敢犯令。没入田童。﹂。\\
\hspace*{2em}天子乃思卜式之言。召拜式为中郎。爵左庶长。赐田十顷。布告天下。使明知之。\\
\hspace*{2em}初。卜式者。河南人也。以田畜为事。亲死。式有少弟。弟壮。式脱身出分。独取畜羊百余。田宅财物尽予弟。式入山牧十余岁。羊致千余头。买田宅。而其弟尽破其业。式辄复分予弟者数矣。是时汉方数使将击匈奴。卜式上书。原输家之半县官助边。天子使使问式。﹁欲官乎。﹂式曰。﹁臣少牧。不习仕宦。不愿也。﹂使问曰。﹁家岂有冤。欲言事乎。﹂式曰。﹁臣生与人无分争。式邑人贫者贷之。不善者教顺之。所居人皆从式。式何故见冤于人。无所欲言也。﹂使者曰。﹁苟如此。子何欲而然。﹂式曰。﹁天子诛匈奴。愚以为贤者宜死节于边。有财者宜输委。如此而匈奴可灭也。﹂使者具其言入以闻。天子以语丞相弘。弘曰。﹁此非人情。不轨之臣。不可以为化而乱法。愿陛下勿许。﹂于是上久不报式。数岁。乃罢式。式归。复田牧。岁余。会军数出。浑邪王等降。县官费众。仓府空。其明年。贫民大徙。皆仰给县官。无以尽赡。卜式持钱二十万予河南守。以给徙民。河南上富人助贫人者籍。天子见卜式名。识之。曰﹁是固前而欲输其家半助边﹂。乃赐式外繇四百人。式又尽复予县官。是时沃尓沃皆争匿财。唯式尤欲输之助费。天子于是以式终长者。故尊显以风百姓。\\
\hspace*{2em}初。式不愿为郎。上曰。﹁吾有羊上林中。欲令子牧之。﹂式乃拜为郎。布衣屩而牧羊。岁余。羊肥息。上过见其羊。善之。式曰。﹁非独羊也。治民亦犹是也。以时起居。恶者辄斥去。毋令败群。﹂上以式为奇。拜为缑氏令试之。缑氏便之。迁为成皋令。将漕最。上以为式朴忠。拜为齐王太傅。\\
\hspace*{2em}而孔仅之使天下铸作器。三年中拜为大农。列于九卿。而桑弘羊为大农丞。筦诸会计事。稍稍置均输以通货物矣。\\
\hspace*{2em}始令吏得入谷补官。郎至六百石。\\
\hspace*{2em}自造白金五铢钱后五岁。赦吏民之坐盗铸金钱死者数十万人。其不发觉相杀者。不可胜计。赦自出者百余万人。然不能半自出。天下大抵无虑皆铸金钱矣。犯者众。吏不能尽诛取。于是遣博士褚大。徐偃等分曹循行郡国。举兼并之徒守相为利者。而御史大夫张汤方隆贵用事。减宣。杜周等为中丞。义纵。尹齐。王温舒等用惨急刻深为九卿。而直指夏兰之属始出矣。\\
\hspace*{2em}而大农颜异诛。初。异为济南亭长。以廉直稍迁至九卿。上与张汤既造白鹿皮币。问异。异曰。﹁今王侯朝贺以苍璧。直数千。而其皮荐反四十万。本末不相称。﹂天子不说。张汤又与异有郄。及有人告异以它议。事下张汤治异。异与客语。客语初令下有不便者。异不应。微反唇。汤奏当异九卿见令不便。不入言而腹诽。论死。自是之后。有腹诽之法比。而公卿大夫多谄谀取容矣。\\
\hspace*{2em}天子既下缗钱令而尊卜式。百姓终莫分财佐县官。于是告缗钱纵矣。\\
\hspace*{2em}郡国多奸铸钱。钱多轻。而公卿请令京师铸钟官赤侧。一当五。赋官用非赤侧不得行。白金稍贱。民不宝用。县官以令禁之。无益。岁余。白金终废不行。\\
\hspace*{2em}是岁也。张汤死而民不思。\\
\hspace*{2em}其后二岁。赤侧钱贱。民巧法用之。不便。又废。于是悉禁郡国无铸钱。专令上林三官铸。钱既多。而令天下非三官钱不得行。诸郡国所前铸钱皆废销之。输其铜三官。而民之铸钱益少。计其费不能相当。唯真工大奸乃盗为之。\\
\hspace*{2em}卜式相齐。而杨可告缗遍天下。中家以上大抵皆遇告。杜周治之。狱少反者。乃分遣御史廷尉正监分曹往。即治郡国缗钱。得民财物以亿计。奴婢以千万数。田大县数百顷。小县百余顷。宅亦如之。于是商贾中家以上大率破。民偷甘食好衣。不事畜藏之产业。而县官有盐铁缗钱之故。用益饶矣。\\
\hspace*{2em}益广关。置左右辅。\\
\hspace*{2em}初。大农筦盐铁官布多。置水衡。欲以主盐铁。及杨可告缗钱。上林财物众。乃令水衡主上林。上林既充满。益广。是时越欲与汉用船战逐。乃大修昆明池。列观环之。治楼船。高十余丈。旗帜加其上。甚壮。于是天子感之。乃作柏梁台。高数十丈。宫室之修。由此日丽。\\
\hspace*{2em}乃分缗钱诸官。而水衡。少府。大农。太仆各置农官。往往即郡县比没入田田之。其没入奴婢。分诸苑养狗马禽兽。及与诸官。诸官益杂置多。徒奴婢众。而下河漕度四百万石。及官自籴乃足。\\
\hspace*{2em}所忠言。﹁世家子弟富人或斗鸡走狗马。弋猎博戏。乱齐民。﹂乃征诸犯令。相引数千人。命曰﹁株送徒﹂。入财者得补郎。郎选衰矣。\\
\hspace*{2em}是时山东被河灾。及岁不登数年。人或相食。方一二千里。天子怜之。诏曰。﹁江南火耕水耨。令饥民得流就食江淮间。欲留。留处。﹂遣使冠盖相属于道。护之。下巴蜀粟以振之。\\
\hspace*{2em}其明年。天子始巡郡国。东度河。河东守不意行至。不辨。自杀。行西逾陇。陇西守以行往卒。天子从官不得食。陇西守自杀。于是上北出萧关。从数万骑。猎新秦中。以勒边兵而归。新秦中或千里无亭徼。于是诛北地太守以下。而令民得畜牧边县。官假马母。三岁而归。及息什一。以除告缗。用充仞新秦中。\\
\hspace*{2em}既得宝鼎。立后土。太一祠。公卿议封禅事。而天下郡国皆豫治道桥。缮故宫。及当驰道县。县治官储。设供具。而望以待幸。\\
\hspace*{2em}其明年。南越反。西羌侵边为桀。于是天子为山东不赡。赦天下囚。因南方楼船卒二十余万人击南越。数万人发三河以西骑击西羌。又数万人度河筑令居。初置张掖。酒泉郡。而上郡。朔方。西河。河西开田官。斥塞卒六十万人戍田之。中国缮道馈粮。远者三千。近者千余里。皆仰给大农。边兵不足。乃发武库工官兵器以赡之。车骑马乏绝。县官钱少。买马难得。乃着令。令封君以下至三百石以上吏。以差出牝马天下亭。亭有畜牸马。岁课息。\\
\hspace*{2em}齐相卜式上书曰。﹁臣闻主忧臣辱。南越反。臣愿父子与齐习船者往死之。﹂天子下诏曰。﹁卜式虽躬耕牧。不以为利。有余辄助县官之用。今天下不幸有急。而式奋愿父子死之。虽未战。可谓义形于内。赐爵关内侯。金六十斤。田十顷。﹂布告天下。天下莫应。列侯以百数。皆莫求从军击羌。越。至酎。少府省金。而列侯坐酎金失侯者百余人。乃拜式为御史大夫。\\
\hspace*{2em}式既在位。见郡国多不便县官作盐铁。铁器苦恶。贾贵。或强令民卖买之。而船有算。商者少。物贵。乃因孔仅言船算事。上由是不悦卜式。\\
\hspace*{2em}汉连兵三岁。诛羌。灭南越。番禺以西至蜀南者置初郡十七。且以其故俗治。毋赋税。南阳。汉中以往郡。各以地比给初郡吏卒奉食币物。传车马被具。而初郡时时小反。杀吏。汉发南方吏卒往诛之。间岁万余人。费皆仰给大农。大农以均输调盐铁助赋。故能赡之。然兵所过县。为以訾给毋乏而已。不敢言擅赋法矣。\\
\hspace*{2em}其明年。元封元年。卜式贬秩为太子太傅。而桑弘羊为治粟都尉。领大农。尽代仅筦天下盐铁。弘羊以诸官各自市。相与争。物故腾跃。而天下赋输或不偿其僦费。乃请置大农部丞数十人。分部主郡国。各往往县置均输盐铁官。令远方各以其物贵时商贾所转贩者为赋。而相灌输。置平准于京师。都受天下委输。召工官治车诸器。皆仰给大农。大农之诸官尽笼天下之货物。贵即卖之。贱则买之。如此。富商大贾无所牟大利。则反本。而万物不得腾踊。故抑天下物。名曰﹁平准﹂。天子以为然。许之。于是天子北至朔方。东到太山。巡海上。并北边以归。所过赏赐。用帛百余万匹。钱金以巨万计。皆取足大农。\\
\hspace*{2em}弘羊又请令吏得入粟补官。及罪人赎罪。令民能入粟甘泉各有差。以复终身。不告缗。他郡各输急处。而诸农各致粟。山东漕益岁六百万石。一岁之中。太仓。甘泉仓满。边余谷诸物均输帛五百万匹。民不益赋而天下用饶。于是弘羊赐爵左庶长。黄金再百斤焉。\\
\hspace*{2em}是岁小旱。上令官求雨。卜式言曰。﹁县官当食租衣税而已。今弘羊令吏坐市列肆。贩物求利。亨弘羊。天乃雨。﹂。\\
\hspace*{2em}太史公曰。农工商交易之路通。而龟贝金钱刀布之币兴焉。所从来久远。自高辛氏之前尚矣。靡得而记云。故书道唐虞之际。诗述殷周之世。安宁则长庠序。先本绌末。以礼义防于利。事变多故而亦反是。是以物盛则衰。时极而转。一质一文。终始之变也。禹贡九州。各因其土地所宜。人民所多少而纳职焉。汤武承毙易变。使民不倦。各兢兢所以为治。而稍陵迟衰微。齐桓公用管仲之谋。通轻重之权。徼山海之业。以朝诸侯。用区区之齐显成霸名。魏用李克。尽地力。为强君。自是以后。天下争于战国。贵诈力而贱仁义。先富有而后推让。故庶人之富者或累巨万。而贫者或不厌糟糠。有国强者或并群小以臣诸侯。而弱国或绝祀而灭世。以至于秦。卒并海内。虞夏之币。金为三品。或黄。或白。或赤。或钱。或布。或刀。或龟贝。及至秦。中一国之币为二等。黄金以溢名。为上币。铜钱识曰半两。重如其文。为下币。而珠玉。龟贝。银锡之属为器饰宝藏。不为币。然各随时而轻重无常。于是外攘夷狄。内兴功业。海内之士力耕不足粮饷。女子纺绩不足衣服。古者尝竭天下之资财以奉其上。犹自以为不足也。无异故云。事势之流。相激使然。曷足怪焉。\\

% ===== 索隐述赞（司马贞《史记索隐》卷二十九）=====
\clearpage
\shuzanhead{索隐述赞曰}\linebreak
{\shuzanbody\fontsize{24}{41}\selectfont
\hbox to \hsize{\hfil 平准之立。通货天下。既入县官。或振华夏。\hfil}\linebreak
\hbox to \hsize{\hfil 其名刀布。其文龙马。增算告缗。裒多益寡。\hfil}\linebreak
\hbox to \hsize{\hfil 弘羊心计。卜式长者。都内克殷。取赡郊野。\hfil}%
\par}

\end{document}
